\documentclass[12pt,a4paper]{article}
\usepackage[utf8]{inputenc}
\usepackage[portuguese]{babel}
\usepackage{hyperref}

\title{Content}
\author{Tomás Pinto - Scripting no Processamento de Linguagem Natural}
\date{\today}

\begin{document}

\maketitle

\begin{abstract}
Com base no modelo de n-grams desenvolvido e no método de scoring aplicado, as três frases selecionadas do texto são:
\begin{itemize}
    \item As  opiniões  sobre  o  FEC  podem,  assim,  mudar  consoante  o  indivíduo  e  a comunidade em  que  este se  insere.
    \item Mas não só os conflitos contribuem para a alteração das fronteiras dos países que compõem o continente europeu.
    \item Apesar de menos comum, é possível encontrarmos este mesmo feito, em 1997, com a Grécia e a Turquia.
\end{itemize}
\end{abstract}

\section{Texto Original com Entidades (NER)}
Introdução O  Festival  Eurovisão  da  \textbf{Canção} (LOC)  é  uma  das  manifestações  artísticas  e  culturais mais populares no espaço europeu e além-fronteiras, alcançando um papel fulcral quando pensamos na cultura e no modo como esta pode ser apresentada àqueles que, por vezes, se encontram mais distantes, em termos geográficos. Esta competição coloca em palco diferentes  povos,  com  tradições  e  costumes  distintos,  unidos  em  prol  da  música  e  da convivência.

É  com  este  propósito  que  desenvolvemos  esta  investigação,  para  tentarmos perceber de que forma é vivido o \textbf{Festival Eurovisão da Canção} (MISC) e de que modo se poderá apreender um sentido de multiculturalidade no continente europeu.  Para  contextualizar este concurso recorreremos a elementos históricos, a fim de perceber como é que o povo europeu  se  sente  em  relação  aos  seus  semelhantes.  Para  tal,  entendemos  necessário conhecer  a  história  do  \textbf{Velho  Continente} (LOC)  para  que  se  possa  compreender  a  realidade cultural refletida nos tempos de hoje. Neste sentido, começaremos esta dissertação pela história do continente europeu, para refletirmos sobre os diversos povos e o seu papel na fundação  de  grande  parte  dos  elementos  e  relações  que  nos  identificam  enquanto indivíduos de uma nação. Assim, tentaremos perceber o contexto em que surgiu o festival e de que forma o passado contribuiu e fundamentou a sua criação. Procuraremos também analisar o modo como o público vive esta competição, para tentarmos perceber se existe um sentimento de unidade e/ou diversidade inerente ao mesmo.

Para que seja possível uma abrangência maior a nível de resultados, utilizaremos diferentes metodologias. A pesquisa bibliográfica, tanto física como digital, está na base da metodologia da nossa investigação. Recorreremos ainda à realização de um inquérito online para  melhor  tentar  perceber  a  opinião  dos  telespectadores  acerca  do  \textbf{Festival Eurovisão  da  Canção} (MISC).  Propomo-nos,  pois,  entender  o  modo  como  uma  manifestação artística congrega em si povos e culturas distintos. Num primeiro momento, partimos de uma análise histórica que permitirá elencar os países e como estes comunicam desde os primórdios até à atualidade, para tentar compreender se o passado em comum dos mesmos facilita a cumplicidade no espaço europeu. Em seguida, daremos uma maior atenção ao \textbf{Festival Eurovisão da Canção} (MISC) e de que forma este promove a cultura, seja ela individual, seja ela europeia. O inquérito que referimos anteriormente irá completar as informações que desenvolveremos. Será que fatores culturais, tais como a língua, podem se apresentar como promotores da individualização de um país? De que forma o passado histórico pode facilitar a manutenção da comunicação ou dificultá-la? São estas e outras questões que nos propomos desenvolver ao longo da nossa investigação.

A história da \textbf{Europa} (LOC) não é linear e também não o é a sua identidade possível.

\textbf{Vasco Graça Moura} (PER) 1. A \textbf{Europa} (LOC) no mundo contemporâneo Quando refletimos sobre o continente europeu, percebemos que a sua história é vasta e que  se desenvolveu ao longo  de vários séculos, tendo-se tornado  na referência cultural  dos  nossos  dias.  Contudo,  se  pretendermos  refletir  sobre  a  sua  fundação,  não podemos  deixar  de  analisar  a  mitologia  grega  e,  por  conseguinte,  o  mito  de  \textbf{Europa} (LOC).

\textbf{Europa} (LOC) era filha de \textbf{Agenor} (PER), rei da \textbf{Fenícia} (LOC), e de \textbf{Teléfassa} (LOC). O mito conta-nos que, perante o encontro de \textbf{Zeus} (PER) com \textbf{Europa} (LOC), numa praia, este fica arrebatado pela sua beleza ao vê- la entre as diferentes pessoas que a rodeavam. Foi nesse preciso momento que, motivado pelo desejo de tê-la consigo, se terá transformado num touro branco. Ao aproximar-se de \textbf{Europa} (LOC),  esta  decide  tocar  nele  e  depois  de  perceber  que  este  não  constituía  qualquer ameaça sobre si, decide subir para o dorso do animal. Com este movimento por parte da jovem, \textbf{Zeus} (PER) desata a correr pelo mar fora até à ilha de \textbf{Creta} (LOC). Aquando da sua chegada, o deus  assume  a  sua  forma  normal.  Ambos  decidem  permanecer  naquele  local,  na  atual \textbf{Grécia} (LOC), continente europeu, deixando as raízes fenícias, ou o que consideraríamos a \textbf{Ásia} (LOC).

Enquanto isto, \textbf{Agenor} (PER) e os seus quatro outros filhos faziam todos os possíveis para poder encontrar  \textbf{Europa} (LOC),  que  havia  desaparecido  misteriosamente  no  oceano.  Apesar  dos esforços,  ninguém  conseguiu  encontrá-la.  \textbf{Sabemos} (MISC)  que  dá  à  luz  três  filhos,  \textbf{Minos} (PER), \textbf{Sarpédon} (PER) e \textbf{Radamante} (PER), graças a essa nova união. Muitos investigadores acreditam que esta princesa está na origem do que hoje consideramos a \textbf{Europa} (LOC). Este mito permite-nos um  retrocesso,  na  medida  em  que,  "remete  para  a  cultura  antiga,  para  o  passado  […] recorda histórias de deuses e de heróis, tem uma  tonalidade nebulosa, lírica,  agressiva […]"  (\textbf{Jabouille} (LOC),  1994:  13).  Devemos  ter  em  conta  também  que  "[o]  mito  pode  ainda assumir  uma  dimensão  histórica:  nele  o  homem  reconhece  as  suas  origens,  as  suas tradições […]" (\textbf{Jabouille} (LOC), 1994: 17) e que, \textbf{"} (PER)[a] mitologia clássica grega e latina continua a ser a grande fonte referencial do nosso universo mítico (…)" (\textbf{Jabouille} (LOC), 1993: 7).

Como  apontado, os  mitos  estão frequentemente  associados à  fundação de algo.

Todavia, deveremos ver em que sociedade se enquadram. Os \textbf{Gregos} (LOC) e os \textbf{Romanos} (LOC) foram dois dos povos mais influentes no continente europeu desde os primórdios dos tempos.

Observavam  a  cultura,  reservando-lhe  um  lugar  de  destaque.  Ambas  as  civilizações podem  ser  consideradas  o  berço  de  diferentes  elementos  que,  nos  tempos  de  hoje,  se destacam como importantes, como é o caso do direito, da política ou da tecnologia. Para a  nossa  investigação,  iremos  debruçar-nos  sobretudo  em  três  aspetos  específicos:  a linguagem, a arte e a religião.

A  civilização  grega  abriu  as  portas  para  a  civilização  romana,  partilhando elementos  comuns ou  elementos  que decorrem  do anterior,  como  é o caso da religião.

Encontramos o culto do fantástico, isto é, os deuses, que se distinguiam entre si, no caso das  culturas  gregas  e  romanas.  A  título  de  exemplo,  podemos  referir  que  na  \textbf{Grécia} (LOC), encontramos \textbf{Zeus} (PER), como o pai dos deuses e dos homens, equivalendo a \textbf{Júpiter} (LOC), em \textbf{Roma} (LOC), e muitos outros casos. Contudo, com o surgimento da religião cristã, que analisaremos mais adiante, no \textbf{Império Romano} (LOC), o \textbf{Cristianismo} (MISC) passou a deter mais importância e a ser um  ponto  fulcral  no  desenvolvimento  da  sociedade  ocidental.  Também  a  linguagem surgiu  com  estes  povos,  sobretudo  as  que  perduram  no  continente  europeu,  apesar  de todas  as  transformações  que  sofreram  desde  esses  tempos,  como  explicaremos  num próximo  ponto.  Nos diferentes  elementos relacionados com  o mundo  artístico, as duas nações  serviram  de  base  para  todos  os  processos  que  estiveram  na  origem  cultural  de diferentes civilizações. O teatro foi, talvez, na história do mundo grego, uma das maiores fontes de manifestações culturais, tendo sido fundamental para outras que se seguiram.

Também  os  diferentes  traços  arquitetónicos,  as  esculturas,  as  danças,  as  pinturas  e  as obras  literárias  contribuíram  para  a  forte  permanência  e  influência  de  certos  traços característicos  destas  sociedades.  \textbf{Percebemos} (LOC)  que  a  cultura  vai  para  além  do  que consideramos  enquanto  movimentos  culturais,  sendo  a  cultura  nestes  períodos  toda  a envolvente de um determinado povo, num certo local.

A religião, como já foi referido, teve um papel fundamental e passou por diferentes fases, sendo os deuses, em certa medida, num primeiro momento objeto de admiração e veneração de um ser superior ao \textbf{Homem} (PER). Para que possa haver uma nova etapa no que toca  à  religião  em  contexto  europeu,  devemos  relembrar  dois  movimentos,  um  destes europeu,  ligado  à  filosofia  grega,  e  outro  proveniente  do  exterior,  a  religião  judaica,  a forma religiosa mais antiga que conseguimos encontrar. Um nome que poderíamos vir a associar é \textbf{Fílon de Alexandria} (PER), um filósofo que uniu estes pensamentos: \textbf{Many} (LOC) of his writings are extant. \textbf{They} (PER) offer a unique access to a mode of thinking that is based on the Bible and \textbf{Jewish} (ORG) traditions, but is deeply influenced by \textbf{Hellenism} (PER) and in particular the doctrines of \textbf{Greek} (MISC) philosophy. His allegorical \textbf{method of exegesis and his} (MISC) theology in turn exercised a strong \textbf{influence on the thought of the early Church Fathers} (MISC), \textbf{who ensured the survival of his works} (ORG) (\textbf{David T.} (PER), 2019: s/p).

A religiosidade passou a deter, com o passar dos anos, uma maior influência na sociedade em que se encontrava contígua. A fé permitiu aos indivíduos uma crença em algo superior, que poderia guiar os mesmos. Foi em resultado do que acabámos que referir que surgiu o \textbf{Cristianismo} (MISC), uma das maiores demonstrações de religiosidade presente no continente europeu até os dias de hoje. Esta forma assenta no culto e no credo na vida de \textbf{Jesus  Cristo} (PER)  e  dos  ensinamentos  que  estão  presentes  no  grande  livro  desta  religião,  a \textbf{Bíblia} (MISC).  No  entanto,  houve  períodos  em  que  esta  religião  viu  o  seu core ser  afetado.  A divisão dos cristãos em católicos, ortodoxos e protestantes, desde o século XIV/XV, veio abrir a possibilidade de fundação de, além de outras formas de crença num ser superior, ramificar uma dada religião consoante o local em que esta mesma é praticada: […] as \textbf{Religiões} (LOC) existem, são um facto incontornável, têm a maior importância no destino das sociedades que as adoptam e são, por outro lado, mais ou menos influenciadas pelas culturas dessas sociedades (\textbf{Carmo} (LOC), 2001: 27).

Em particular, no século \textbf{XXI} (MISC), percebemos que a religião perspetivou e conseguiu, mesmo que sem  esse propósito,  aumentar  esses  ramos de que falamos, estando alguns deles ligados ao extremismo, especialmente em regiões como o \textbf{Médio Oriente} (LOC), \textbf{África} (LOC), entre  outros.  Essas  regiões  contêm  diferentes  visões  religiosas,  contribuindo  para  uma confluência e, por vezes, choque de perspetivas diferentes.

Ao  longo  de  todos  os  séculos,  é  possível  encontrarmos  períodos  conturbados ligados  a  disputas,  guerras  ou  pandemias,  que  foram  prejudiciais  para  o  continente europeu.  Disso  pode  ser  exemplo  a  \textbf{Revolução  Francesa} (MISC),  em  pleno  século  XVIII,  pelo impacto  que  teve,  não  só  na  nação  francesa,  como  também  a  nível  mundial.  Esta revolução assentava na ideia de liberdade extensível a todos os indivíduos, sem exceção, do povo francês, relativamente a amarras invisíveis que ditavam e limitavam as escolhas individuais.  A  monarquia  absolutista  foi  derrubada,  afastando  os  privilégios  feudais, aristocráticos e religiosos, e assim abrindo caminho para uma governação nacionalista, invocando  os  ideais  iluministas  de  \textbf{Liberdade} (LOC),  Igualdade  e  \textbf{Fraternidade} (LOC).  Ainda  assim, foram  diversos  os  conflitos  armados  que  esta  potência  travou  contra  algumas  pátrias europeias que ainda se encontravam submetidas à monarquia. Exemplo de alguns desses países podem ser \textbf{Áustria} (LOC), \textbf{Portugal} (LOC), \textbf{Rússia} (LOC), \textbf{Espanha} (LOC), \textbf{Países Baixos} (LOC), \textbf{Reino Unido} (LOC), entre outros.  A  história  política  do  \textbf{Velho  Continente} (LOC)  viria  a  mudar  drasticamente  após  a \textbf{Revolução  Francesa} (MISC),  demonstrando  que  poderia  haver  outros  meios,  afastados  dos poderes  tradicionais,  para  uma  soberania  política,  mais  justa  e  nacionalista.  Mais recentemente, e seguindo esta ótica, podemos ainda referir as \textbf{Guerras Mundiais} (MISC), conflitos que marcaram o século XX e que tiveram enormes repercussões no continente europeu.

A par da \textbf{Revolução Francesa} (MISC), vieram a envolver dois grandes blocos, a \textbf{Tríplice Entente} (MISC) e a \textbf{Tríplice Aliança} (MISC) que dividiram o continente europeu em dois, tanto na \textbf{Primeira} (MISC) como na  \textbf{Segunda  Grande  Guerra} (MISC).  Estes  blocos  ainda  hoje  mantêm  alianças  entre  si  e apresentam um grande poder de decisão em alguns órgãos que foram criados após esse conflito, com o objetivo de manter a paz internacional.

Para  que  fosse  possível  um  armistício,  como  expusemos  no  parágrafo  anterior, houve  a  necessidade  da  criação  de  instituições  europeias  e  internacionais  que contribuíssem para a garantia da democracia, da igualdade, do respeito e da solidariedade entre os membros. Foi assim que, em 1951, nasceu a \textbf{Comunidade Europeia do Carvão e do Aço} (ORG) (\textbf{CECA} (ORG)), designada, posteriormente de \textbf{União Europeia} (ORG) (\textbf{UE} (ORG)).  Esta organização apresenta um grau crescente no que toca à sua importância para os europeus, estando na mira  de  diferentes  países  que  ainda  não  fazem  parte  dessa  comunidade  e  que  desejam poder vir a integrar.

A \textbf{Europa} (LOC) dos dias de hoje concentra em si uma enorme variedade multicultural ligada  aos  povos  que  a  compõem  e  aos  que  estiveram  na  sua  fundação  enquanto continente. Neste sentido, os processos de desenvolvimento e fixação das suas fronteiras nacionais parecem-nos relevantes para que possamos analisar o seu contributo e impacto num determinado país ou região. É possível afirmar que as nações europeias que integram a \textbf{Europa} (LOC) deste século contêm vestígios de elementos distintos deixados pelos povos que habitaram primordialmente este continente, como já tivemos oportunidade de explicar.

Estes  elementos  podem  estar  ligados  ao  passado  histórico,  tanto  do  local  como  da população,  as  tradições  culturais  que  perduraram  e  foram  transmitidas  de  geração  em geração,  a  própria  escrita  e  língua  de  um  país  e  a  sua  religião.  É  claro  que,  com  o desenvolvimento  da  sociedade,  ao  longo  dos  diferentes  momentos  que  abordámos, podemos  aprofundar  as  semelhanças,  mas  também  as  diferenças  entre  as  nações.  O caminho  para  a  unidade  intracontinental  irá  residir  não  só  nas  diferenças  e  nas semelhanças,  mas  também  na  forma  como  estas  nações  cooperam  entre  si  para  a manutenção  da  paz  e  como  se  aproximam  ou  afastam  consoante  os  elementos  que elencámos anteriormente. Estes pontos, aliados ao \textbf{FEC} (ORG), poderão ajudar a demonstrar de que forma a cultura de cada país contribui para a multiculturalidade europeia, questão que constitui  objeto  da  nossa  atenção  e  à  qual  tentaremos  dar  resposta  ao  longo  da  nossa investigação.

1.1.

\textbf{Fronteiras Geográficas} (LOC) Quando  observamos  um  mapa  geográfico  do  continente  europeu  é  possível constatar que, na sua constituição, ao longo dos diferentes séculos, os países integrantes e as suas fronteiras foram-se alterando pelas mais variadas razões. Nos dias de hoje, se pretendermos  delimitar  qual  é  o  território  considerado  europeu  iremos  ter,  a  norte,  o \textbf{Oceano  Glacial  Ártico} (LOC),  a  sul,  o  \textbf{Mar  Mediterrâneo} (LOC)  e  o  \textbf{Mar  Negro} (LOC),  a  oeste,  o  \textbf{Oceano Atlântico} (LOC), e, a leste, os  \textbf{Montes Urais} (LOC),  localizados na \textbf{Rússia} (LOC), havendo uma parte russa europeia e uma parte russa asiática.

Contudo, nem sempre foi assim tão claro, a nível territorial, delimitar a extensão desta região devido aos impérios que se consolidaram ao longo da história. Encontramos no  século  \textbf{V  a.C.} (MISC),  \textbf{Heródoto} (PER),  historiógrafo  grego,  que,  segundo  \textbf{Cícero} (PER),  seria  o  pai  da \textbf{História  como  a  conhecemos} (MISC).  Baseado  nos  seus  escritos  da  época,  encontramos  uma representação do mundo e, consequentemente, delimitação dos continentes distinta das demais devido aos seus conhecimentos sobre os diferentes povos de cada continente.

Com  o  desenvolvimento  dos  povos,  as  conquistas  e  perdas  territoriais,  o continente  europeu  foi-se  constituindo  com  gentes  provenientes  de  outros  lugares  e portadores de outras culturas. Como abordado anteriormente, os diferentes impérios que ocupavam lugar no que consideramos hoje em dia o continente europeu tiveram um papel fulcral  no  estabelecimento  de  diferentes  elementos,  como  a  língua,  tema  que  iremos desenvolver num ponto seguinte.

Quando refletimos sobre este tópico, afigura-se-nos importante referir três grandes impérios:  \textbf{Romano} (LOC),  \textbf{Otomano} (LOC)  e  \textbf{Califado  Omíada} (MISC).  Tanto  o  \textbf{Império  Romano} (MISC)  como  o \textbf{Império Otomano} (LOC) são-nos mais presentes a nível histórico, uma vez que fazem parte da fundação  de  grande  parte  do  espaço  que  conhecemos  enquanto  \textbf{Europa} (LOC).  O  \textbf{Império Otomano} (LOC),  apesar  de  ter  surgido  séculos  mais  tarde,  é  considerado  fundamental  na preservação e consolidação de diversas tradições das regiões que conquistaram ao longo de cerca de três séculos. Os turcos passaram a olhar para os elementos culturais dos locais que  conquistavam  e  traduziram  esses  pontos  numa  nova  forma  cultural  que  passaria  a definir o que é a identidade cultural desse império. Apesar da adição de novos elementos fruto deste movimento, podemos encontrar ainda, dos povos primordiais europeus, como os persas, elementos característicos. É na arquitetura que conseguimos observar melhor a influência dos persas e da parte mais oriental europeia ligada ao islão. As mesquitas, os túmulos e outros pontos culturais detinham lembranças de componentes característicos de outras partes europeias como o \textbf{Barroco} (LOC), com forte presença na arquitetura portuguesa, ou o \textbf{Rococó} (LOC), ligado à cultura francesa, sendo \textbf{"} (ORG)como [um] desdobramento do barroco, mais  leve  e  intimista  que  aquele  e  foi  usado  inicialmente  em  decoração  de  interiores" (\textbf{Imbrosi} (PER),  2023:  s/p).  Já  o  \textbf{Califado  Omíada} (MISC)  apresenta-se  como  um  império  menos abordado na história europeia. Isto pode suceder-se devido a duas razões: por se localizar maioritariamente em áreas mais orientais do continente europeu e/ou por estar mais ligado à religião islâmica: É  indubitável  que  os  homens  elaboraram  culturas  diferentes  em  virtude  do  seu afastamento geográfico, das propriedades particulares do meio e da ignorância em que se encontravam  em  relação  ao  resto  da  humanidade,  mas  isso  só  seria  rigorosamente verdadeiro se cada cultura ou cada sociedade estivesse ligada e se tivesse desenvolvido no isolamento de todas as outras (\textbf{Lévi-Strauss} (PER), 1996: 15/16).

Todos estes processos tornaram o \textbf{Velho Continente} (LOC) num local de confluência de diferentes culturas e tradições. De certo modo, encontramos uma aprimoração do conceito que  atualmente  é  muito  utilizado,  multiculturalidade.  O  multiculturalismo  está diretamente associado à globalização e ao fenómeno de êxodo de um determinado local.

"Multiculturalismo  não  designa  apenas  uma  realidade  fixa  e  homogénea,  clara  e univocamente  configurável  por  alguns  traços  essenciais  susceptíveis  de  serem universalizados a todos os posicionamentos que enfrentam e que o tematizam" (\textbf{André} (LOC), 2012: 26). O multiculturalismo favorece um melting-pot , ou seja, através da convivência entre diferentes gentes, com passados e tradições distintas, poderemos vir a descobrir uma nova  identidade  que  pode  estar  ligada  a  um  determinado  local,  como  aconteceu  nos impérios  que  anteriormente  listámos.  A  aceitação  de  tradições  e  elementos  culturais exteriores a um determinado povo permite-lhes a convivência com os povos conquistados de uma maneira mais simples e contribui para promover as diferenças. Contudo, devemos compreender que este conceito pode concentrar em si uma pluralidade de significados, tais como uma aceção descritiva, ligada à dimensão descritiva de multiculturalismo, e a aceção normativa, provocando controvérsia: Entretanto, o contexto em que convocámos para esta reflexão o fenómeno da globalização é o de uma reflexão sobre identidade, identidades e multiculturalidade e o que é certo é que, apesar da sua pré-história e da sua história, as características com que tal fenómeno hoje se apresenta, articuladas com o fenómeno da sociedade-em-rede, e as implicações que  tem  sobre  a  esfera  cultural,  marcam  com  novos  traços  o  fenómeno  da multiculturalidade e contribuem também para sublinhar as insuficiências do conceito de cultura a que acabámos de fazer referência (\textbf{André} (LOC), 2012: 39/40).

A  identidade  tal  como  nos  é  apresentada  por  \textbf{José  Maria  André} (PER),  em \textbf{Multiculturalidade} (LOC),  identidades  e  mestiçagem:  o  diálogo  intercultural  nas  ideias,  na política, nas artes e na religião , é algo que é inerente ao indivíduo e que  se apresenta fulcral quando pensamos nas singularidades que cada detém num local de confluência.

Por sua vez,  essa interação entre gentes  com  culturas distintas  promove um  fenómeno denominado  de  mestiçagem,  levando  à  origem  de  seres  mestiços  por  natureza.  Este fenómeno  expandiu-se  largamente  na  época  do  colonialismo,  pelos  europeus,  que chegavam às regiões até então desconhecidas. Em contexto europeu, podemos refletir no conceito de mestiçagem se relembrarmos os primeiros povos colonizadores deste mesmo continente.  Incluindo  nós,  os  portugueses,  podemos  ser  considerados  mestiços, remontando  à  época  das  disputas  pela  \textbf{Península  Ibérica} (LOC),  existindo  ainda,  hoje  em  dia, diversos aspetos relacionados com este mesmo fenómeno, como é o caso de elementos ligados à história árabe e a presença africana, incluindo o \textbf{Fado} (MISC), melodia tão característica da identidade portuguesa. Poucos são os povos que se podem definir como "puros", no sentido de se manterem fiéis a toda a história e ADN do seu povo, existindo somente em regiões mais remotas do globo. As fronteiras físicas de países ou regiões tornam-se mais ténues quando pensamos nas identidades e nas caracterizações de cada povo, embora estas tenham sido revistas ao longo dos séculos devido a disputas que podiam alterar a nossa visão e identificação de cada local: \textbf{The  very} (MISC)  concepts  of  homogenous  national  cultures,  the  consensual  or  contiguous transmission of historical traditions, or 'organic' ethnic communities - as the grounds of cultural comparativism - are in a profound \textbf{process of redefinition} (MISC) (\textbf{Bhabha} (MISC), 1994: 5).

Após a queda dos grandes impérios, o estabelecimento das fronteiras sofreu alguns percalços  ao  longo  dos  tempos.  Se  pensarmos  em  pleno  século  XV,  o layout europeu passou a incluir muitos mais países, frutos das desintegrações imperiais, e já encontramos estabelecidos  alguns  países  que  reconhecemos  atualmente  como  \textbf{Portugal} (LOC),  \textbf{Inglaterra} (LOC), \textbf{Hungria} (LOC) e \textbf{França} (LOC), por exemplo. É de salientar que estes dois últimos países, apesar de os reconhecermos,  não  continham  ainda  as  fronteiras  ou  a  localização,  tal  como  as conhecemos neste século.

\textbf{Importa} (LOC) refletir sobre um dos séculos recentes mais importantes a nível histórico para  a  nossa  sociedade  que  é  o  século  XX.  Cada  nação  experienciou  diferentes acontecimentos, mas, algo em que quase todos participaram foi as \textbf{Guerras Mundiais} (MISC), seja a \textbf{Primeira} (MISC) ou a \textbf{Segunda} (MISC). O clima de instabilidade bélico entre 1914 e 1945, na \textbf{Europa} (LOC), moldou grande parte dos povos que se envolveram nesses conflitos à escala mundial. Foi em resultado da \textbf{Primeira Grande Guerra} (MISC) que as fronteiras europeias se alteraram de forma significativa,  acabando  com  o  reinado  de  grandes  impérios  como  o  alemão,  o  austro- húngaro, o russo e o otomano, mas, permitiu o estabelecimento de novas nações como a \textbf{Áustria} (LOC), a \textbf{Hungria} (LOC), e, inclusive, a independência de muitos estados como a \textbf{Finlândia} (LOC), a \textbf{Polónia} (LOC),  entre  outros.  Foi  graças  às  ligações  que  estes  recém-independentes  estados detinham com os impérios anteriores que moldaram as suas culturas, partindo assim, de um ponto base comum.

A formação da \textbf{Sociedade das Nações} (ORG) (SN), em 1919, possibilitou um começo no tratamento  da  estabilidade  geopolítica  europeia  após  o  conflito  armado.  Contudo,  o ambiente de paz que havia sido produzido pela assinatura dos diferentes tratados de paz foi dissolvido, havendo o despoletar da \textbf{Segunda Guerra Mundial} (MISC), anos mais tarde. Esse acontecimento  é  considerado,  no  campo  histórico,  como  um  dos  mais  marcantes  pelo número de perdas humanas que veio a causar em ambos os lados que se opunham. Por uma segunda vez, no século XX, as fronteiras e os poderes nacionais viram outros chegar ao  seu  poder,  como  foi  o  caso  da  divisão  da  \textbf{Alemanha} (LOC),  um  dos  protagonistas  neste conflito  armado,  em  \textbf{República  Federal  da} (LOC)\textbf{Alemanha} (LOC)  e  \textbf{República  Democrática  da Alemanha} (LOC). Da mesma forma, como vimos a criação da SN após o desfecho da \textbf{Primeira Grande Guerra} (MISC), duas novas alianças militares internacionais foram criadas para colmatar as diferenças e pôr em prática a paz na \textbf{Europa} (LOC), sendo essas a \textbf{Organização do Tratado do Atlântico Norte} (ORG) (\textbf{OTAN} (ORG)), em 1949, e o \textbf{Pacto de Varsóvia} (ORG), em 1955. Apesar dos objetivos comuns para a manutenção de um clima de estabilidade, não só europeu, mas global, a \textbf{Guerra Fria} (MISC) veio colocar tensão nas relações políticas.

Mas não só os conflitos contribuem para a alteração das fronteiras dos países que compõem o continente europeu. Três grandes dissoluções geopolíticas, no fim do século XX, permitiram a independência de distintos lugares. A dissolução da \textbf{União Soviética} (ORG) deu  início  a  um  processo,  em  dezembro  de  1991,  num  clima  de  paz  entre  a  \textbf{Rússia} (LOC),  a \textbf{Ucrânia} (LOC)  e  a  \textbf{Bielorrússia} (LOC).  No  ano  seguinte,  a  \textbf{Checoslováquia} (LOC)  dividiu  o  seu  espaço  de influência em dois países diferentes, a \textbf{República Checa} (LOC) e a \textbf{Eslováquia} (LOC), tendo sido uma divisão em que a equidade reinou, seja a nível de bens do \textbf{Estado} (LOC), seja a nível de outros elementos.  Este  acontecimento  ficou  conhecido  como  o  "\textbf{Divórcio  de  Veludo"} (MISC).  Por último, mas não menos importante para a nossa análise, foi a desintegração da \textbf{Jugoslávia} (ORG), uma região báltica, facto que veio a concentrar em si mais países. Enquanto no caso da \textbf{Checoslováquia} (LOC) o mundo assistiu a uma separação amigável, onde ambas as nações viram o seu poderio ser dividido equitativamente, na \textbf{Jugoslávia} (ORG) deu-se uma série de conflitos e instabilidade,  relacionados  com  a  economia,  problemas  estruturais,  entre  outros,  que levou  esta  dissolução  e  consequente  formação  de  novos  países,  sendo  eles  \textbf{Bósnia  e Herzegovina} (LOC),  \textbf{Croácia} (LOC),  \textbf{Montenegro} (LOC),  \textbf{República  da  Macedónia} (LOC)  (que  viria  a  tornar-se  a \textbf{Macedónia do Norte} (LOC) após conflitos com uma região grega), \textbf{Sérvia} (LOC) e \textbf{Eslovénia} (LOC). Um dos casos que requer uma análise mais atenta prende-se com a situação do \textbf{Kosovo} (LOC), que por alguns países, como \textbf{Portugal} (LOC), é considerado país, e por outros, como \textbf{Espanha} (LOC), não. Todas estas  unificações  e  desintegrações  que  ocorreram  no  continente  europeu  ao  longo  do século  XX  e,  sobretudo,  nas  décadas  finais  deste  mesmo  período,  vieram  a  alterar  a geografia  e  dinâmicas  europeias,  facto  que  irá  influenciar  o  nosso  objeto  de  estudo,  o \textbf{Festival Eurovisão da Canção} (MISC) (\textbf{FEC} (ORG)): E, de facto, a seguir à \textbf{Segunda Guerra Mundial} (MISC) desencadeou-se um intenso movimento de reflexão sobre o espírito europeu, a construção europeia, o humanismo, a violência, a guerra, a cultura, o progresso técnico e o progresso moral (\textbf{Moura} (LOC), 2013: 9).

\textbf{Capaz} (PER) de integrar as fronteiras invisíveis que compõem e sustentam o continente europeu, encontramos as \textbf{Organizações} (LOC), algumas que integram apenas países europeus e, noutros  casos,  extraeuropeus.  \textbf{A  UE} (MISC),  fundada  em  1951,  inicialmente  designada  por \textbf{CECA} (ORG), como referimos no ponto anterior, tem como missão promover a paz em território europeu,  fomentar  a  coesão  territorial,  no  que  concerne  a  nível  económico  e  social, através de um  mercado  único sem  fronteiras. O  seu crescimento  ao longo das décadas levou a que muitos países que integram esta organização tivessem um crescimento seguro em áreas distintas. O respeito e a aceitação de cada indivíduo, isto é, "respeitar a riqueza da diversidade cultural e linguística da  \textbf{UE} (ORG)" (\textbf{União Europeia} (ORG), s/d: s/p), são alguns dos valores  promovidos  pela  mesma.  Com  o  despoletar  da  guerra  na  \textbf{Ucrânia} (LOC),  passámos  a encontrar a \textbf{OTAN} (ORG), ou a \textbf{Aliança Atlântica} (ORG), no centro das negociações no que toca a uma aliança militar e política, promovendo a segurança de todos os seus estados-membros. É com este objetivo que a \textbf{OTAN} (ORG) deve a sua fundação após a \textbf{Segunda Guerra Mundial} (MISC) e congrega países europeus e extraeuropeus. Apesar disso, o conflito bélico recente entre a \textbf{Rússia} (LOC) e a \textbf{Ucrânia} (LOC) veio  despoletar uma necessidade, por parte de determinados países que ainda não faziam parte dessa organização, como a \textbf{Suécia} (LOC) e a \textbf{Finlândia} (LOC), de passar a integrar os mesmos, visto o perigo iminente de conflitos armados.

Outros  exemplos  de  organizações  igualmente  importantes  para  a  convivência intraeuropeia e extraeuropeia podem ser dados através do \textbf{Fundo Monetário Internacional} (ORG) (\textbf{FMI} (ORG)),  a  \textbf{Organização  das  Nações  Unidas} (ORG)(\textbf{ONU} (ORG)),  \textbf{Organização  para  a  Cooperação  e Desenvolvimento Económico} (ORG) (\textbf{OCDE} (ORG)) e a \textbf{Organização Mundial do Comércio} (ORG) (\textbf{OMC} (ORG)).

Estas organizações possibilitam a interligação com gentes de pontos diferentes do globo, e permitem que exista uma interação e confluência de ideias e ideais que podem promover a aceitação do outro de uma forma mais simplificada.

1.2. \textbf{Fronteiras Culturais} (LOC) Como procurámos demonstrar no ponto anterior, é possível encontrar fronteiras geográficas que podem ditar a proximidade de determinados povos. Mas, devemos refletir também  sobre  o  facto  de  que  as  fronteiras  culturais,  mesmo  dentro  de  uma  nação, poderem influenciar os indivíduos. Normalmente, este tipo de delimitação cultural não se observa a olho nu, mas apenas quando começamos a analisar os intervenientes na ação.

É  por  essa  razão  que,  para  a  nossa  investigação,  as  fronteiras  culturais  se  apresentam como  mais  relevantes  face  às  geográficas.  Em  virtude  disso,  existem  elementos fundamentais para a investigação em curso de que podem ser exemplo a língua, a escrita, a religião e as tradições. Porém, devemos refletir sobre o vocábulo que congrega todos estes itens, respetivamente a cultura. Quando procuramos definir o que é a cultura em si ou a cultura de um povo percebemos que este termo pode apresentar em si uma metáfora para todo o processo cultural que ocorre com os seres humanos. O que significa que, mais do que a primeira acepção do conceito cultura associado à \textbf{"} (ORG)acção, processo ou efeito de cultivar\textbf{"} (MISC) (\textbf{Houaiss} (MISC), 2003: 1152), normalmente associado ao cultivo da terra, a cultura vai mais além, podendo também ser definida como um: conjunto  de  padrões  de comportamento, crenças,  conhecimentos,  costumes  […]  (uma) forma ou etapa evolutiva das tradições e valores intelectuais, morais, espirituais (de um determinado lugar ou período) […] (um) complexo de actividades, instituições, padrões sociais ligados à criação e difusão das belas-artes, ciências humanas e afins […] (\textbf{Houaiss} (MISC), 2003: 1152).

Também, \textbf{Mário Vargas Losa} (PER) em \textbf{A Civilização do Espetáculo} (MISC) considera que: Ao longo da história, a noção de cultura teve diferentes significados e matizes. Durante muitos séculos foi um conceito inseparável da religião e do conhecimento teológico; na \textbf{Grécia} (LOC) foi marcado pela filosofia e em \textbf{Roma} (LOC) pelo direito, enquanto no \textbf{Renascimento} (LOC) era impregnado  sobretudo  pela  literatura  e  pelas  artes.  Em  épocas  mais  recentes  como  o iluminismo foram a ciência e as grandes descobertas científicas que deram a orientação principal  à  ideia  de  cultura.  Mas,  apesar  dessas  variantes  e  até  à  nossa  época,  cultura sempre  significou  um  conjunto  de  fatores  e  disciplinas  que,  segundo  amplo  consenso social, a constituíam e ela implicava: a reivindicação de um património de ideias, valores e  obras  de  arte,  de  conhecimentos  históricos,  religiosos,  filosóficos  e  científicos  em constante evolução de novas formas artísticas e literárias e da investigação em todos os campos do saber (\textbf{Losa} (PER), 2012: 61).

Por sua vez, para \textbf{Isabel Ferin} (PER): Na \textbf{Antiguidade Clássica} (PER), o conceito de cultura designa a acção que o homem realiza  - quer sobre o seu meio, quer sobre si mesmo - no sentido de aperfeiçoar as suas qualidades e promover a cultura do espírito (\textbf{Ferin} (LOC), 2009: 35).

A cultura pode ser experimentação e reflexão, pensamento e sonho, paixão e poesia e uma revisão crítica constante e profunda de todas as certezas, convicções, teorias e crenças (\textbf{Losa} (LOC), 2012: 70).

\textbf{Denys  Cuche} (LOC),  antropólogo  e  professor  na  \textbf{Universidade  Paris} (LOC)  \textbf{Descartes} (PER), doutorado em  etnografia, considera na sua obra \textbf{A Noção de Cultura que "} (MISC)A noção de cultura é inerente à reflexão das ciências sociais […] pensarem a unidade na diversidade da humanidade de outro modo que não em termos biológicos" (\textbf{Cuche} (PER), 2006: 23).

Acrescenta ainda que: Ao  contrário  da  noção,  mais  ou  menos  rival  dentro  do  mesmo  campo  semântico,  de sociedade, a noção de cultura aplica-se apenas ao que é humano. E oferece a possibilidade de concebermos a unidade do homem na diversidade dos seus modos de vida e de crenças, incidindo  a  tónica,  conforme  os  investigadores,  ora  mais  na  unidade,  ora  mais  na diversidade (\textbf{Cuche} (LOC), 2006: 25).

\textbf{Refere} (LOC)  as  ideias  alemãs  e  francesas  como  forma  de  compreensão  de  conceções distintas  como  modelos  de  cultura,  associadas  a  conceitos  de  nação,  nacionalismo  e  a conquistas de diversos níveis: A ideia alemã de cultura evolui, portanto, em certa medida, ao longo do século XIX, sob a influência do nacionalismo. Liga-se cada vez mais ao conceito de "nação\textbf{"} (MISC). A cultura releva da alma, do génio de um povo. A nação cultural precede e solicita a nação política.

A cultura surge como como um conjunto de conquistas artísticas, intelectuais e morais que constituem o património de uma nação, considerado como adquirido de uma vez por todas, e que fundam a sua unidade (\textbf{Cuche} (LOC), 2006: 37) O debate franco-alemão entre os séculos XVIII e XIX é arquetípico das duas concepções de cultura, uma particularista e a outra universalista, que se encontram na base das duas maneiras de definir o conceito de cultura nas ciências sociais contemporâneas (\textbf{Cuche} (LOC), 2006: 38).

\textbf{Isabel  Ferin} (PER),  que  apresentámos  anteriormente,  demonstra-nos  que    \textbf{"} (PER)[…]  o conceito de cultura passa a estar associado ao conceito de civilização, opondo-se-lhe ou complementando-o,  mas  sempre  enfatizando  o  movimento  em  direcção  ao desenvolvimento do indivíduo e das sociedades\textbf{"} (MISC) (\textbf{Ferin} (MISC), 2009: 36). \textbf{Refere} (LOC) também, em \textbf{Comunicação  e} (MISC)\textbf{Culturas} (LOC)  do  \textbf{Quotidiano} (LOC), de  que  forma  \textbf{Malinowski} (PER),  um  antropólogo polaco ligado à fundação da escola funcionalista e ao estudo da etnografia, define cultura.

\textbf{Malinowski} (PER) define cultura como o todo integral que compõem os instrumentos e os bens de consumo, as castas constitutivas dos vários reagrupamentos sociais, as ideias, as artes, crenças  e  costumes  atribuindo-lhe  a  função  de  satisfazer  necessidades  de  carácter biológico/natural ou simbólicas (\textbf{Ferin} (MISC), 2009: 38).

De modo sintético, face aos diferentes sociólogos contemporâneos, como \textbf{Talcott Parsons} (PER), \textbf{Ferin} (MISC) refere que: a sua concepção da \textbf{Teoria da Cultura} (MISC), ao longo de quatro décadas de produção, \textbf{Talcott Parsons} (PER)  definirá  genericamente  cultura  como  um  sistema  complexo  e  relativamente coerente de significados, normas e valores que orientam a acção social (\textbf{Ferin} (LOC), 2009: 43).

Se associarmos as diferentes definições do conceito de cultura, apresentadas por \textbf{Isabel  Ferin} (PER),  \textbf{Denys} (LOC)  Cuche  ou  \textbf{Mário  Vargas  Losa} (PER),  podemos  afirmar  que  a  cultura  se reflete  como  uma  aprendizagem,  desde  o  berço  para  quem  nasceu  dentro  de  uma determinada  cultura,  como  é  o  caso  da  maioria  dos  indivíduos  ou,  uma  aprendizagem quando passam a integrar outra cultura que não aquela em que nasceram, como acontece com os emigrantes. Apesar desta diferença, compreendemos que por se apresentar como uma  caminhada  que  envolve  diferentes  ramos,  pode  apresentar  dificuldades  para  o indivíduo.  Neste  sentido,  devemos  repensar  estes  mesmos  ramos  para  que  possamos encontrar,  na  nossa  investigação,  elementos  que  contribuam  para  um  sentimento  de unidade e/ou diversidade, dentro do continente europeu. As etapas do conceito de cultura associadas à história dos povos e dos países proporcionam o conhecimento e o progresso quando esta procura adaptar-se à atualidade.

O primeiro sinal que encontramos quando nos deslocamos a um determinado lugar e que se pode apresentar como uma fronteira é a língua. Esta permite-nos identificar um povo, podendo conter em si diferentes dialetos ou variantes, mesmo em locais vizinhos.

Por vezes, conseguimos denotar que \textbf{"} (MISC)é raro que os limites das línguas naturais coincidam com  as  fronteiras  políticas  dos  vários  \textbf{Estados} (LOC)  onde  são  faladas"  (\textbf{Walter} (PER),  1996:  273).

Trata-se de  um  "sistema de representação constituído  por palavras e por  regras que  as combinam  em  frases  que  os  indivíduos  de  uma  comunidade  linguística  usam  como principal meio de comunicação e de expressão, falado ou escrito\textbf{"} (MISC) (\textbf{Houaiss} (MISC), 2003: 2283).

Diferentes  línguas  ou  formas  de  escrita  perduram  no  \textbf{Velho} (LOC)Continente  desde  os  seus primórdios, estando na génese das formas de expressão que vieram desaguar no século \textbf{XXI} (MISC). Todas estas formas linguísticas compõem a família das línguas indo-europeias. Se quisermos distinguir linguisticamente as nações que integram a \textbf{Europa} (LOC), encontraremos três grupos distintos: as germânicas, as românicas e as eslavas. Se pensarmos nas origens remotas que estas mesmas têm somos levados, respetivamente, às \textbf{Invasões Germânicas} (MISC), ao \textbf{Império Romano} (LOC) e à \textbf{Expansão Celta} (MISC). Apesar de não integrar estes conjuntos, sabemos que  a  civilização  grega  e,  por  conseguinte,  o  dialeto  grego,  teve  um  papel  central  na constituição  da  sociedade  pelo  que  "foi  o  berço  da  civilização  ocidental,  [que] desempenhou um papel essencial na implantação e difusão do alfabeto , palavra que, de resto, testemunha bem a sua origem grega ( alfa-beta ) na estrutura etimológica que possui" (\textbf{Walter} (LOC),  1996:  31).  As  línguas  célticas,  não  obstante,  na  sua  representatividade  no continente europeu não se encontram a par e passo com as românicas e as germânicas, concentrando-se, na sua maior parte, em territórios regidos pela coroa britânica como a \textbf{Ilha de Man} (LOC), a \textbf{Escócia} (LOC) ou a \textbf{Península da Cornualha} (LOC), entre outros. Os romanos seguiram as bases deixadas pelos gregos e deram origem a uma das línguas base da comunicação, o latim, que \textbf{"} (MISC)se tornou através dos séculos, sob a sua forma escrita e erudita, a língua por excelência da cultural ocidental" (\textbf{Walter} (LOC), 1996: 95).

Nada à partida, contudo, faria prever tão excelente destino para a língua deste pequeno povo constituído por aldeões agricultores cujos povoamentos humildes não passavam, em meados do séc. \textbf{VIII a.C.} (PER), de um modesto lugar de passagem no vale pantanoso do \textbf{Tibre} (LOC), situado no coração do \textbf{Lácio} (LOC) (\textbf{Walter} (LOC), 1996: 95).

Apesar  de  ser  considerada  uma  língua  morta,  o  latim  engloba  diferentes metamorfoses  que  o  levaram  a  estar  na  origem  de  grande  parte  das  línguas  europeias, como é o caso do português, do italiano, do romeno ou do castelhano. O romeno é uma língua que, apesar das "inúmeras invasões eslavas dos séculos VI e \textbf{VII} (PER) […] mantém uma base latina, apesar dos seus inúmeros traços eslavos" (\textbf{Walter} (LOC), 1996: 121). A língua latina foi fortemente influenciada em diversos aspetos, seja através de outras línguas, como é o caso da germânica, que iremos abordar mais adiante, ou de povos, podendo destacar os \textbf{Bizantinos} (MISC)  ou  os  \textbf{Árabes} (LOC),  respetivamente  " metro (medida  -  do  grego \textbf{metron} (PER) , medida)"  (\textbf{Walter} (PER),  1996:  127)  e  " ribes (groselha,  embora  a  palavra  designasse originariamente o ruibarbo)" (\textbf{Walter} (LOC), 1996: 128).

No que diz respeito às línguas germânicas é possível identificar períodos distintos na sua história. \textbf{Percebemos} (LOC) que houve uma ligação intrínseca entre os povos germânicos e  os  outros  povos  que  já  abordamos  a  nível  da  sua  língua.  Essa  troca  permitiu  que houvesse  alguns  elementos  comuns  com  cada  tipo  linguístico,  como  por  exemplo,  a \textbf{"} (ORG)palavra alemã \textbf{Tisch} (LOC) (mesa) foi colhida do latim \textbf{DISCUS} (PER) (tabuleiro), o que remete para o costume romano de servir aos convivas tabuleiros já guarnecidos que se retiravam no final das refeições" (\textbf{Walter} (LOC), 1996: 276). Também o grego se encontra refletido nesta conexão indiretamente, na medida em que, "o termo latino MONACHUS (monge), por exemplo,  proveniente  do  termo  grego monachos (inicialmente  solitário,  depois monge)  que  esteve  na  origem  do  topónimo \textbf{München} (LOC) (\textbf{Munique} (LOC))"  (\textbf{Walter} (PER),  1996: 276). Mas esta influência funciona para ambos os intervenientes, sendo possível encontrar no latim algumas réstias ligadas ao germânico, \textbf{"} (PER)como é o caso de \textbf{RENO} (LOC) (pele de rena), \textbf{SAPO} (MISC) (sabão), \textbf{GANTA} (ORG) (gansa) e \textbf{GLAESUM} (MISC) (âmbar)" (\textbf{Walter} (LOC), 1996: 276).

Nos nossos dias, o dinamarquês, o neerlandês, o alemão e o finlandês são algumas das  línguas  oficiais  que  provêm  desta  descendência.  Podemos,  assim,  identificar  no continente  europeu  determinadas  localizações  que  estão  associadas  a  um  ramo linguístico. No extremo ocidental encontramos o ramo céltico, no norte europeu o ramo germânico, mais a sul o ramo itálico/românico e a sudeste o ramo helénico. É percetível que,  apesar  das  distâncias  entre  as  nações,  já  desde  os  tempos  antes  de  \textbf{Cristo} (PER)  até  aos nossos dias, existem elos no que diz respeito às línguas e às influências que cada um teve nas outras, independentemente do contexto em que estes ocorreram. A herança linguística permite  que,  apesar  de  conterem  palavras  distintas  a  nível  escrito,  que  exista  uma perceção auditiva de cada elemento, estejam ou não dentro da mesma família linguística.

Se colocarmos este ponto em análise, refletindo em povos afastados, do ponto de vista geográfico, mas que fazem parte da mesma família a nível da língua, sendo eles \textbf{Portugal} (LOC), \textbf{Espanha} (LOC) e \textbf{França} (LOC) face  à \textbf{Roménia} (LOC), encontramos  palavras semelhantes, respetivamente, jornalista - \textbf{jurnalist} (PER), \textbf{autobus} (PER) - autobus e fenêtre - fereastră.

A  língua  é  fundamental  para  a  singularidade  de  uma  população  e  é  através  da língua que podemos distinguir, logo à partida, uma participação dentro do nosso objeto de  estudo,  o  \textbf{FEC} (MISC).  \textbf{Acresce} (LOC)  também  a  escrita  que  está  inteiramente  ligada  à  nação  e  é através da escrita que os povos se mantêm na história de um lugar. As pinturas rupestres deram, em certa medida, início a este processo, passando de desenhos para um alfabeto desenhado com runas, para um alfabeto com grafemas semelhantes aos que encontramos nos dias de hoje, mas característicos dos povos primordiais. A escrita é o \textbf{"} (ORG)acto ou efeito de  escrever  ou  de  redigir  (…)  por  meio  de  signos  gráficos"  (\textbf{Houaiss} (MISC),  2003:  1569), permitindo  aos  povos  ter  na  sua  posse  um  documento  que  prova  a  sua  presença  num contexto histórico, seja esse ligado ao continente europeu, seja num âmbito mais global.

A  religião,  como  já  apontamos  anteriormente,  apresenta-se  como  um  ponto relevante e que contribui para a identidade específica de uma nação ou, dependendo do seu  espaço  de  influência,  região.  "Não  há  a  certeza  de  que  religio  apresente  o  mesmo radical  de  religare,  conforme  se  comenta  no \textbf{Dicionário  Houaiss} (MISC) ,  em  nota  etimológica baseada no \textbf{Dictionnaire Etymologique de la Langue Latine} (MISC) , de \textbf{A. Ernout} (PER) e A. Meillet:1 \&\#34;o prefixo é re-, red- (cf. relliquiae, \textbf{reliquiae)\&\#34} (MISC);, dizem \textbf{Ernout} (MISC) e \textbf{Meillet} (MISC), \&\#34;mas o segundo elemento é obscuro. Os latinos ligam-no a relegere (...), etimologia defendida por \textbf{Cícero} (PER) (...). Outros autores [\textbf{Lactâncio} (PER) e \textbf{Sérvio} (MISC)] associam religĭo a religāre : seria propriamente `o fato de se ligar com relação aos deuses, simbolizado pela utilização das uittae [`fitas para enfeitar as vítimas ou ornar os \textbf{altares} (LOC)] e dos stémmata no culto. Alega-se em favor desse sentido a imagem de \textbf{Lucrécio} (PER), 1, 931: religionum nodis animum exsoluere ; (...). O sentido seria portanto: `obrigação assumida para com a divindade; vínculo ou escrúpulo religioso (cf. mihi religio est `tenho o escrúpulo de); depois `culto prestado aos deuses, religião" (\textbf{Rocha} (LOC), 2012: s/p). A religião \textbf{Católica} (ORG) prevalece em grande parte dos países europeus.  Contudo,  "(…)  the  English,  \textbf{German} (PER),  \textbf{Swiss} (LOC),  \textbf{Hungarian} (LOC),  and  \textbf{Netherlandic} (PER) groups  are  noteworthy  for  the  coexistence  of  \textbf{Roman  Catholicism  and  Protestantism} (PER)" (\textbf{Britannica} (MISC), s/d: s/p). Apesar disso, o \textbf{Judaísmo} (MISC) e o \textbf{Islamismo} (MISC), num raio mais pequeno de ação, tiveram um papel de destaque no \textbf{Velho Continente} (LOC). Um dos casos mais famosos destas religiões pertence à \textbf{Segunda Guerra Mundial} (MISC) que, em certa medida, opunha a raça ariana, isto é, os designados "alemães puros\textbf{"} (MISC), e os judeus, considerados por \textbf{Hitler} (PER) como "vermes  parasitas  que  deveriam  ser  eliminados"  (\textbf{United  States  Holocaust  Memorial Museum} (ORG),  s/d:  s/p).  Porém,  esta  visão  sobre  este  povo  não  estava  ligada  diretamente  à religião,  que  seria  considerada  anti  judaica,  embora  fosse  antissemitismo.  O  que diferencia estes dois factos é o ódio desta religião face à sua ideologia e discriminação perante os indivíduos como se estes fossem menores que os seus opressores.

No  entanto,  para  o  elemento  em  estudo,  o  \textbf{FEC} (MISC),  a  religião  não  ocupa  um  dos lugares de destaque no que concerne à distinção ou preferência por uma performer versus outra.  As  diferentes  religiões  são  maioritariamente  vistas  quando  pensamos  num determinado  país  ou  na  sua  localização,  uma  vez  que  afirmámos  que  países  do  norte europeu  encontram  as  duas  formas  de  religião  ligadas  à  religião  romana  ou  cristã, enquanto no sul se destaca o catolicismo clássico, com \textbf{Roma} (LOC) como centro. Não obstante, a tolerância religiosa em alguns pontos do continente torna-se cada vez mais importante, para que exista um crescente sentimento de pertença a um lugar. A iniciativa mais recente para levar este propósito mais além foi a do governo britânico que, "pela primeira vez, \textbf{Londres} (LOC) [celebrou] o mês sagrado do \textbf{Ramadão} (LOC) com iluminação pública" (\textbf{Sic Notícias} (MISC), 2023: s/p).

Da  mesma  forma  que  a  língua  nos  permite  identificar  uma  cultura  ou  um determinado  povo  no  momento  em  que  o  observamos,  o  vestuário  desempenha igualmente um papel semelhante, visto que, \textbf{"} (PER)compõem o traje\textbf{"} (MISC) (\textbf{Houaiss} (MISC), 2003: 3696) que poderá estar associado a uma respetiva nação e à sua identidade. Desde sempre, as vestimentas tiveram um carácter de distinção entre os indivíduos, seja pela sua qualidade, quando  pensamos  nos  reis  e  o  seu  povo,  seja  pelas  suas  tonalidades.  A  vestiaria  dos europeus está composta por cores como o vermelho, que simboliza o \textbf{"} (ORG)princípio da vida, com a sua força, o seu poder e o seu brilho, o vermelho, cor de fogo e de sangue […] ambivalência de onde provem todo o poder de fascínio da cor vermelha, que traz em si intimamente ligadas as duas mais profundas pulsões humanas: acção e paixão, libertação e  opressão"  (\textbf{Chevalier} (PER),  1994:  686/687),  e  o  dourado,  associado  ao  ouro  que,  segundo \textbf{"} (PER)dizem os brâmanes, é a imortalidade […] [uma] identificação com a luz solar, o ouro sempre foi um dos símbolos de \textbf{Jesus} (PER), \textbf{Luz} (LOC), \textbf{Sol} (LOC), \textbf{Oriente} (MISC)" (\textbf{Chevalier} (PER), 1994: 495/497).

Da mesma maneira que conseguimos encontrar semelhanças linguísticas devido à sua história no \textbf{Velho Continente} (LOC), existem elementos em comum nos trajes típicos, não só específicos de um país, mas também intraeuropeus. As afinidades podem encontrar-se a diferentes níveis, seja na coloração dos fatos, na sua estética, sendo um vestido ou um conjunto entre uma saia e uma camisa, ou umas calças e uma camisa, ou, inclusive, numa \textbf{Portugal} (LOC) e \textbf{República Checa} (LOC).

parecença  a  nível  dos  acessórios  que  utilizam,  podendo  incluir  objetos  de  outros  em materiais  como  o  ouro  em  algumas  regiões  mais  ligadas  ao  leste  europeu.  Abaixo, apresentamos  alguns  exemplos  de  indumentárias  europeias  de  diferentes  pontos  para melhor tentarmos ilustrar este tema.

Todavia,  encontramos,  de  igual  modo,  pontos  distintos  que  seguem  a  ideia  de pensamento que demonstramos acima.

1 \textbf{Países Baixos} (LOC) e \textbf{Suécia} (LOC).

As semelhanças ou as diferenças que conseguimos encontrar nos trajes típicos não se  apresentam  como  algo  negativo.  \textbf{Proporcionam} (PER),  por  sua  vez,  uma  diferenciação positiva, na maioria dos casos. Sendo o continente europeu diversificado nas suas gentes e  culturas,  é  normal  encontrarmos  pessoas  com  indumentárias  características  de  uma nação ou de uma determinada religião. Cada vez mais, graças às migrações populacionais, a  \textbf{Europa} (LOC)  encontra-se  embebida  nessa  mesma  realidade  o  que,  por  vezes,  se  pode apresentar como um problema no que tange a discriminação devido a eventos anteriores que ocorram inclusive em solo europeu.

\textbf{Algo} (LOC) que também podemos encontrar como base de uma determinada cultura está relacionado com as tradições que estas contêm. Há que salientar que algumas tradições podem  estar  diretamente  associadas  ao  passado  histórico  ou  a  uma  religião.  Podemos encontrar uma "comunicação oral de factos, lendas, ritos, usos, costumes etc. de geração para geração\textbf{"} (MISC) (\textbf{Houaiss} (MISC), 2003: 3555) que, através do seu carácter repetitivo ao longo dos anos,  se  torna  num  costume,  passando  a  ser  característico  e  identificativo  de  uma população. O elemento  abordado  anteriormente, o vestuário,  funciona como  elemento- chave para ligarmos certas tradições a um local e a um povo. Fora as tradições típicas, ligadas ao culto religioso, se se tratar de um país em que a religião arrecada um papel importante,  as  celebrações  ou  festas  populares  que  se  relacionam  com  um  elemento identificativo de um \textbf{Estado} (LOC) ganham proporções desmedidas, levando à própria recriação das mesmas em outros pontos do globo. Um facto que poderá encontrar a sua justificação desde  a  segunda  metade  do  século  XIX,  por  termos  encontrado  diversos booms migratórios de fora da \textbf{Europa} (LOC) para dentro e vice-versa. Um dos casos mais evidentes é o \textbf{Dia de São Patrício} (MISC) ou \textbf{St} (MISC). \textbf{Patrick's} (PER) Day , com raízes irlandesas e base religiosa. Apesar de se tratar de uma festividade que serve sobretudo para festejar a comunidade e a cultura irlandesa  marca  também  a  chegada  da  religião  \textbf{Cristã} (MISC)  com  esta  entidade,  \textbf{São  Patrício} (LOC).

Com  a  colonização  da  \textbf{América} (LOC),  por  parte  dos  britânicos,  houve  uma  migração  de tradições que foram deixadas aos colonizados. Sendo assim, \textbf{The earliest recorded} (MISC) parade was held in 1601 \textbf{in what is now St} (MISC). \textbf{Augustine} (LOC), \textbf{Florida} (LOC). The parade, and a \textbf{St} (MISC). Patrick's Day celebration a year earlier, \textbf{were organized by the Spanish Colony's Irish} (MISC) vicar \textbf{Ricardo Artur} (PER) (\textbf{History.com Editors} (MISC), 2017: s/p).

O boom da  migração  de  \textbf{Irlandeses} (LOC)  no  século  XIX  para  os  \textbf{Estados  Unidos  da América} (LOC) esteve envolto em discriminação, o que passou a conferir a estas manifestações culturais ligadas a \textbf{São Patrício} (PER) um momento de exaltação da cultura e das tradições que lhes haviam sido incutidas enquanto ainda residiam na \textbf{Irlanda} (LOC).

As  tradições  estão  na  base  da  manutenção  da  identidade  cultural  e  assim  da comunidade e são essas que têm o poder de unificação quando existe uma deslocação.

Por exemplo, uma comunidade portuguesa, que resida no estrangeiro, mas que mantenha grande parte dos costumes que lhes foram incutidos quando ainda residiam \textbf{em Portugal} (LOC) ou pelos pais, que são portugueses, irá permitir ao indivíduo um sentimento de pertença mais  facilitado  por  rever  os  seus  hábitos  nesta  nova  coletividade  em  que  se  encontra integrado.  Mais  adiante,  iremos  analisar  outros  fenómenos  semelhantes  a  este  mesmo exemplo.

Como  pudemos  observar  anteriormente,  podemo-nos  deparar  com  elementos culturais que se apresentam semelhantes ou com a mesma fundação em partes distintas do \textbf{Velho Continente} (LOC). São estas semelhanças que contribuem para efeito de unidade entre as populações, fomentando um clima de paz, estabilidade e de complementaridade entre os vizinhos, sejam efetivos ou vizinhos por outra distinção. Mas não só as semelhanças têm o poder de aproximar ou distanciar um povo, já que as diferenças arrecadam um papel vital quando refletimos sobre isso. Muitas vezes, é através das diferenças ou da perceção da  visão  do  outro  que  conseguimos,  de  certo  modo,  aproximarmo-nos  dos  demais.  As instituições que abordámos no ponto anterior, tais como a \textbf{UE} (ORG), \textbf{ONU} (ORG), \textbf{OTAN} (ORG), entre outras, contêm  países  com  história,  fundações  e  princípios  distintos,  mas  juntam-se  com  um único propósito que está relacionado com o bem comum. Em situações de crise, é através da observação das diretrizes emanadas por diferentes chefes de \textbf{Estado} (LOC) e  detentores do poder  que  encontramos  pontos  para  melhorar  algum  acontecimento.  Se  recordarmos  o início  da  pandemia  \textbf{COVID-19} (LOC),  as  tomadas  de  decisão  relativamente  à  saúde  pública ditaram  o modo  como  esta doença viria a impactar essa  localidade. Um dos casos  que pode servir como exemplo é o da \textbf{Itália} (LOC) que, com consequente fecho dos estabelecimentos, levou  a  que  muitos  considerassem  que  não  seria  preocupante,  o  que  promoveu  o exponencial desenvolvimento de casos de infeção. Países como \textbf{Portugal} (LOC) e países do norte europeu, como a \textbf{Finlândia} (LOC), a \textbf{Noruega} (LOC) e a \textbf{Suécia} (LOC), destacaram-se por uma ação rápida que fez  com  que  a  propagação  do  vírus  fosse  mais  tardia  e  atenuada.  As  diferentes mentalidades de uma determinada cultura podem promover certos comportamentos que, no ponto de vista de outros, podem à partida não se enquadrar como correto.

Desde 2001 que, na sua generalidade, as migrações e os povos distintos dos que residem num determinado local estão no centro das discussões quando nos lembramos de situações que foram catastróficas e com perdas de vidas humanas devido a atentados. A título  de  exemplo,  podemos  referir  o  11  de  setembro  de  2001,  nos  \textbf{Estados  Unidos  da América} (LOC).  A  visão  que  muitos  detinham  de  elementos  de  outras  religiões  ou  nações apresentava-se  como  um  entrave  para  a  sociedade.  Apesar  de  distante,  também  no continente  europeu,  este  sentimento  de  instabilidade  e  de  receio  da  possibilidade  de ocorrer algo semelhante no futuro se atravessou na mente dos europeus, causando assim, diversos distúrbios e casos de xenofobia perante os outros. A 13 de novembro de 2015, catorze anos após a queda das \textbf{Torres Gémeas} (LOC), \textbf{Paris} (LOC) foi tomada por atentados terroristas ao longo de toda a cidade, sendo no teatro pariense, o \textbf{Bataclan} (MISC), o local que registou o maior número de perda de vidas humanas. Este foi considerado um dos maiores atentados, na  última  década,  em  \textbf{França} (LOC),  embora,  desde  2015,  se  tenha  tornado  mais  recorrente encontrarmos situações neste mesmo contexto e nesse mesmo país. Muitos consideram que este caso está ligado à aceitação de pessoas não francesas que passaram a residir nesse país. Inclusive, em \textbf{Paris} (LOC), conseguimos encontrar numa determinada zona, ruas inteiras ocupadas por pessoas que chegaram a \textbf{França} (LOC) como migrantes. Independentemente disso, um pouco por todo o território europeu, é possível observar ataques que poderiam colocar em causa a segurança e a estabilidade que se havia consolidado após a \textbf{Segunda Guerra Mundial} (MISC). Países como \textbf{Bélgica} (LOC), \textbf{Alemanha} (LOC), \textbf{Espanha} (LOC) e \textbf{Finlândia} (LOC)  também viram no solo nacional  ocorrer  situações  semelhantes.  \textbf{A  UE} (MISC)  afirma  que  é  necessário  "intensificar esforços  para  criar  uma política  europeia  de  migração  que  seja  eficaz,  humanitária  e segura"  (\textbf{Conselho  Europeu} (ORG),  s/d:  s/p).  Contudo,  na  ótica  de  alguns  países,  que  não  se sentem confortáveis a aceitar migrantes, o conflito interior pode surgir e levar a haver a necessidade de outro país, onde as quotas de aceitação já tinham sido ultrapassadas, vir a aceitar mais migrantes. Estes fogem com o pouco que têm, seja da \textbf{guerra da Síria} (MISC) ou do \textbf{Afeganistão} (LOC) ou de outras zonas de conflito, e arriscam-se a atravessar o \textbf{mar Mediterrâneo} (LOC) para  poder  dar  às  suas  famílias  um  local  mais  seguro  para  a  vida.  Esta  crise  tomou proporções globais, tornando-se num dos temas mais debatidos na última década. "Mais de 1.200 pessoas morreram, em 2022, a cruzar o \textbf{Mediterrâneo} (LOC) numa tentativa de chegar à \textbf{Europa} (LOC)" ( \textbf{Público} (LOC) , 2023: s/p), sendo o caso de uma menina síria, de quatro anos, que esteve à deriva vários dias e que acabou perdendo a vida. Muitos pais enviam os filhos, na esperança de um dia poderem reencontrar-se com os mesmos.

Porém,  quando  existe  uma  imigração  noutros  moldes,  podem  surgir  outras questões. Se refletimos sobre a condição da \textbf{Mulher no Médio Oriente} (MISC), na medida em que na sua cultura existe uma subjugação desta perante o homem, semelhante à definição de cidadão que encontrávamos na \textbf{Roma Antiga} (LOC), o pensamento europeu distancia-se dessa forma, não se revendo na inferiorização da mulher. Também a sua religião se apresenta relevante para a visão que o homem tem sobre o sexo feminino visto muitas das práticas se encontrarem descritas no \textbf{Alcorão} (MISC).

As fronteiras culturais podem, por vezes, representar uma dificuldade, mas, por outras  perspetivas,  também  uma  aproximação  pelas  diferenças.  Em  certas  situações, como  foi  o  despoletar  da  guerra  entre  a  \textbf{Rússia} (LOC)  e  a  \textbf{Ucrânia} (LOC),  percebemos  que  os elementos  culturais,  sobretudo  destes  dois  países  que  são  muito  semelhantes  na  sua história,  vieram  a  estar,  em  pleno  século  \textbf{XXI} (LOC),  em  pontos  opostos.  Apesar  disso, devemos  observar  estas  fronteiras  culturais  como  um  propulsor  para  a multiculturalidade  europeia  e  para  o  crescimento  interpessoal  diverso  e  cooperativo entre os diferentes integrantes.

A \textbf{Europa} (LOC) […] será um laboratório com condições ideais para esse modelo, para a allge- meine \textbf{Vereinigung der Menschheit} (MISC) [unificação universal da \textbf{Humanidade} (ORG)] de \textbf{Kant} (PER). A \textbf{Europa} (LOC), com algumas limitações, terá essa capacidade de não só viver como prosperar, com a diversidade cultural (\textbf{Himmel} (PER), 2017: 26).

A  presença  da  multiculturalidade  que  procuramos  abordar  ao  longo  desta investigação,  por  residir  no  seio  da  cultura  europeia  enquanto  unidade,  leva-nos  a repensar em todos os seus movimentos de exôdo e de que forma a comunicação evoluiu ao longo dos diferentes séculos. A diversidade cultural no interior do continente europeu está  em  contante  mudança,  não  se  podendo  caracterizar  como  "uma  realidade  fixa  e homogénea, clara e univocamente configurável por alguns traços essenciais susceptíveis de serem universalizados a todos os posicionamentos que enfrentam e que o tematizam" (\textbf{André} (LOC),  2012:  26).    Em  certa  medida,  esta  multiculturalidade  e  o  facto  de  nos  ser possível  observar  aspetos  distintos  entre  os  indivíduos  num  raio  mais  pequeno  pode permitir um  elogio, isto é, uma exaltação das diferenças, no sentido de aproximar os mesmos pelos aspetos distintos que cada um encontra em si. Como se pode ler no \textbf{Artigo 1} (MISC) da \textbf{Declaração Universal da Diversidade Cultural: Cultural} (MISC) diversity: the common heritage of humanity \textbf{Culture} (MISC) takes diverse forms across time and space. \textbf{This diversity is embodied in the uniqueness and plurality of the identities of the groups and societies making} (MISC) up humankind. As a source of exchange, \textbf{innovation and creativity} (MISC), cultural diversity is as necessary for humankind as biodiversity is for nature. \textbf{In} (MISC) this sense, \textbf{it is the common heritage of humanity and should be recognized and affirmed for the benefit of present and future} (MISC) generations (\textbf{Stenou} (PER), 2002: 4).

Contudo, é possível observar diferentes fenómenos que, através da globalização de um determinado elemento, tendem a standardizar diversos aspetos que englobam a vida de um indivíduo. A título de exemplo podemos referir os estrangeirismos. É prática comum o indivíduo passar a adaptar o seu vocabulário em função de alguns vocábulos que  não  pertencem  à  sua  língua,  mas  que  são  reconhecidos  em  múltiplos  níveis, deixando  de  os  traduzir  para  a  língua  em  que  o  mesmo  passa  a  ser  utilizado.  A comunicação adapta-se aos seus falantes, facilitando o diálogo cultural entre os mesmos pelo reconhecimento das expressões mais identificativas. Como afirma \textbf{Isabel Ferin} (PER): Comunicação como a \textbf{Cultura} (MISC) [integram] sistemas (conjuntos de elementos, objectos ou entidades  que  se  inter-relacionam  mutuamente  para  formas  um  todo  único)  altamente complexos e interdependentes, levando a que qualquer alteração num dos sectores gere alterações, mudanças, rupturas ou adaptações em todos os outros (\textbf{Ferin} (MISC), 2009: 48).

Como  tivemos  oportunidade  de  abordar  anteriormente,  o  europeu  reconhece, pelos seus exemplos históricos, de que forma os diferentes passados culturais podem contribuir para uma não manutenção da paz quando o enfoque reside nas diferenças dos mesmos.  Apesar  disto,  e  pelos  fenómenos  de  conquista  além-fronteiras  e  territórios europeus,  como  no  continente  americano,  africano  ou  asiático,  é  necessário encontrarmos  a  visão  por  pontos  semelhantes  para  que  se  forme  um  diálogo  para  a compreensão e não uma tomada de posição perante os diferentes povos. É certo que, se pretendermos educar e perpetuar os nossos ideais e costumes numa civilização que não os conhece, o diálogo deve assemelhar-se nos seus pontos históricos, mas, de modo a não  errar  com  os  mesmos,  assemelhar  alguns  aspetos  que  estes  pudessem  ter concentrados  nas  suas  culturas  regionais.  Este  fenómeno  pode  ser  repercutido,  por exemplo, através do plano gastronómico. Existem regiões que foram conquistadas pelos europeus, onde a sua história, cultura e sociedade foram estudadas, mas, ao retornarem ao  seu  local  de  origem,  levam  frutos  desta  conquista  que  contribuem  para  que  estes indivíduos se sintam aceites por aqueles que os conquistaram.

\textbf{Song in the history of the European nation is neither} (MISC) simply an object nor a \textbf{subject given meaning} (MISC) through collective performance. Song mobilizes nationalism in exceptionally complex forms, enacting \textbf{the performance of the nation in the ordinary and the extraordinary moments of history} (MISC).

\textbf{Philip Vilas Bohlman} (PER) 2.

Festival Eurovisão da Canção O  Festival  Eurovisão  da  Canção  (\textbf{FEC} (ORG))  surgiu  num  contexto  de  pós-guerra, embrenhado nas mudanças e ideias que o povo europeu pretendia colocar em prática com vista  a  um  sentimento  de  união  em  contexto  \textbf{intracontinental} (LOC).  Neste  sentido,  foi influenciado  pelas  políticas  unionistas  e  pró-europeias  dos  anos  50,  do  século  XX.

\textbf{Importa} (LOC) salientar que a história europeia e a história do \textbf{FEC} (ORG) caminham lado a lado por diferentes razões. Possibilitaram-nos um olhar distinto sobre de que forma uma produção radiotelevisiva  poderia  contribuir  para  a  convivência  e  cordialidade  intraeuropeia.

Apresenta-se, assim, como uma outra forma para a manutenção da paz, numa determinada medida, visto que a \textbf{"} (LOC)[…] cultura é - ou era, quando existia - um denominador comum, algo que mantinha viva a comunicação entre pessoas muito diversas […]" (\textbf{Losa} (PER), 2012: 66).  A  comunicação  entre  os  diferentes  povos  permite  o  diálogo,  abrindo-se  à compreensão.  Apesar  disso,  a  manutenção  da  paz  que  surgiu  com  o  fim  da  \textbf{Segunda Guerra  Mundial} (MISC)  opõe-se  às  situações  críticas  que  começaram  a  despoletar  no  fim  do século XX e início do século \textbf{XXI: Even at the historical moments when it would be obvious to draw attention to cultural diversity and political transformation} (MISC), not least \textbf{during the transition of socialist states in} (MISC) \textbf{Eastern Europe} (ORG) after 1989, \textbf{it is more striking how few} (MISC), \textbf{rather than how many} (MISC), national  entries  have  responded  to  the  pressing  \textbf{issues  of  the  day} (MISC)[…]  (\textbf{Bohlman} (PER), 2007: 47).

Podemos designar por festival um evento que seja "relativo a ou próprio de festa […] em que há alegria, prazer; aprazível, afável, jucundo\textbf{"} (MISC) (\textbf{Houaiss} (MISC), 2003: 1728). Neste sentido,  podemos  refletir  sobre  o  \textbf{FEC} (MISC)  enquanto  um  exemplo  de  festa  que  pretende elucidar distintas pessoas no que toca a diferentes culturas, dos diferentes países europeus, e não só. Este surgiu após uma época conturbada, a \textbf{Segunda Guerra Mundial} (MISC), vivida por grande parte dos países europeus. Mesmo após a sua fundação, a \textbf{Guerra Fria} (MISC), anos mais tarde, contribuiu para um clima de instabilidade. Num período em que todo o continente se encontrava disperso, este concurso possibilitou unir os diferentes povos com um único propósito: apresentar a sua cultura e as suas gentes, para que assim pudessem encontrar nas  suas  diferenças  alguns  pontos  em  comum,  levando  à  cooperação  entre  as  distintas nações  europeias.  Foi  com  este  objetivo  que  surgiu  a  \textbf{European  Broadcasting  Union} (MISC) (\textbf{EBU} (MISC))  em  1950.  Este  serviço  público  de  \textbf{Rádio  Transmissão  Europeia} (LOC)  permitiu  uma maior  proximidade  aos  diferentes  estados  europeus.  Na  sua  fundação,  é  possível encontrarmos alguns dos primeiros integrantes do \textbf{FEC} (ORG) como, por exemplo, a \textbf{Bélgica} (LOC), a \textbf{Itália} (LOC)  ou  a  \textbf{Suíça} (LOC).  Contudo,  antes  de  \textbf{EBU} (MISC),  tínhamos  o  seu  antecessor,  a  \textbf{International Broadcasting  Union} (ORG)  (\textbf{IBU} (MISC)).  Esta  consequente  transformação  deu-se  devido  à  \textbf{Segunda Guerra Mundial} (MISC). A anterior administração encontrava-se insatisfeita com todo o contexto que experienciavam que, apesar do crescente número de rádios ao longo do continente europeu, não foi suficiente para ultrapassar este entrave.

\textbf{Existiu} (LOC) anteriormente uma organização com um objetivo semelhante a ambas, a \textbf{International Radio and Television Network} (ORG) (\textbf{Eastern Europe} (ORG)) (\textbf{ORIT} (ORG)). Nascida na \textbf{Europa} (LOC) de Leste, criada na mesma altura que a \textbf{IBU} (MISC), continha muitos dos países que transitaram, após  o  seu  fecho,  para  a  \textbf{EBU} (MISC).  Nesse  sentido,  com  a  transformação  da  \textbf{IBU} (MISC)  em  \textbf{EBU} (MISC), assistimos ao (re)nascer de um \textbf{Media} (LOC) sem barreiras transfronteiriças, que veio permitir um  maior  alcance.  No  leque  de  diferentes  programas  que  viriam  a  apresentar, encontramos  o  evento  que  constitui  objeto  da  nossa  dissertação:  o  \textbf{FEC} (MISC).  Este  nome, lançado pelo jornalista, \textbf{George Campey} (PER), tornou-se na imagem de marca  do que temos em  voga  nesta  iniciativa  musical.  Se  analisarmos  os  componentes  da  sua  designação, \textbf{"Eurovisão} (MISC)",  encontramos  dois  termos,  "euro"  e  "visão".  Como  uma  possível interpretação, esta designação inclui o objetivo de ver ou observar algo de fonte europeia.

Esta foi a core finalidade da competição, colocando nas radiotransmissoras nacionais, de países distintos, um pouco de cada cultura ou nação. Caso esta denominação não surgisse, o que atualmente é conhecido como \textbf{FEC} (PER), seria \textbf{Festival de Canções de Trocas Televisivas Continentais} (PER), do inglês, \textbf{Continental Television} (ORG) Exchange \textbf{Song Contest} (LOC) : Eurovision  was  a  technical  \textbf{system  of  media  content} (MISC)  interconnections  initially  across \textbf{Europe} (ORG), but later progressively \textbf{linking to the world} (MISC). It allowed content to be \textbf{exchanged} (MISC), \textbf{seen and heard} (MISC), as \textbf{it happened or recorded for later use.} (MISC) The \textbf{Eurovision giant spider's web of connections has grown with demand - reaching out to many others} (MISC), not just \textbf{EBU Members} (MISC). \textbf{In} (MISC) fact, \textbf{the Eurovision technical network} (ORG) now operates as a subsidiary of the EBU - \textbf{Eurovision Services} (MISC) - to serve companies throughout the media and event industry (\textbf{EBU} (MISC), s/d: s/p).

As  opiniões  sobre  o  \textbf{FEC} (MISC)  podem,  assim,  mudar  consoante  o  indivíduo  e  a comunidade em  que  este se  insere. A visão de povos exteriores  ao continente europeu pode  fazer  com  que  esta  manifestação  cultural  não  tenha  tanta  relevância  quanto  em território europeu: \textbf{The Eurovision Song Contest} (ORG) (esc) is not just a pop song competition: often, \textbf{it is} (MISC) a musical manifestation of political and social issues of great significance. […]  \textbf{Although it is often derided  in  popular  discourse -- or  as  in  the} (MISC)United  States,  completely  ignored -- \textbf{Eurovision} (PER) holds a special place in the history of musical nationalism (\textbf{Scott} (PER), 2011: s/p).

Mas são essas mesmas diferentes opiniões que elencam a relevância deste festival enquanto promotor da diversidade cultural. A diversidade de pontos de vista fomenta a opinião pública, permitindo o debate sobre esta manifestação cultural com as proporções mundiais  que  detém  o  \textbf{FEC} (MISC).  Estando  este  festival  no  centro  das  conversas  enquanto movimento  que  promove  cada  nação  enquanto  individual,  ou  através  das  semelhanças que esta consegue encontrar entre as demais, poderemos aumentar o raio de influência que  encontramos  neste  concurso,  levando-o  mais  longe  e  a  públicos  que  poderia  não cativar,  especialmente  através das  novas tecnologias, já que  estas  estão  presentes  cada vez mais no dia a dia de cada indivíduo. Como a opinião pública nos parece fundamental para melhor entendermos o alcance do \textbf{FEC} (MISC), utilizaremos um inquérito que seguirá essa linha  de  pensamento.  Sendo  assim,  analisaremos  os  resultados  advindos  desse  mesmo inquérito mais adiante, aliado a outros tópicos fulcrais para a compreensão do impacto da \textbf{Eurovisão} (LOC).

2.1.

Das origens à atualidade \textbf{A ideia} (MISC) em torno do \textbf{"universo} (MISC)" eurovisivo relaciona-se com o alcance da cultura e com a valorização que é dada a cada elemento cultural. Todos os povos  são diferentes entre si, sendo os indivíduos que os compõem ainda mais distintos, mas é precisamente aí  que  reside  a  nossa  questão  principal  de  investigação:  de  que  forma  a  cultura  pode unificar, através das suas diferenças, um determinado continente?

The social articulation of difference, from the minority perspective, is a complex, on- going negotiation that seeks to authorize cultural hybridities that emerge in moments of historical transformation (\textbf{Bhabha} (MISC), 1994: 2).

Quando refletimos sobre a origem desta competição, podemos ser levados a um país e a um festival específicos. As raízes do \textbf{FEC} (ORG) relacionam-se com o \textbf{Festival de Sanremo} (PER), em  \textbf{Itália} (LOC), visto este ter  servido de inspiração para o formato que temos atualmente. O \textbf{Festival de Sanremo} (PER), que ocorre na comuna de \textbf{Sanremo} (LOC), \textbf{Itália} (LOC), começou em 1951, num país  que  ainda  vivia  com  parcos  recursos  tecnológicos,  como  a  televisão,  mas  que,  ao longo dos anos, passou a ganhar mais interesse aos olhos dos italianos e da população em geral. Tal como o contexto em que surgiu a \textbf{Eurovisão} (LOC), em 1956, o \textbf{Festival de Sanremo} (PER) emergiu num ambiente de instabilidade, apesar de ter decorrido passado cinco anos após o final da \textbf{Segunda Guerra Mundial} (MISC). Uma guerra que levou a que os mais diversos países tivessem sentido repercussões aos mais diferentes níveis, sejam eles económicos, sociais, ambientais  e,  inclusivamente,  culturais.  De  modo  a  tentar  reverter  essa  situação,  a administração da comuna de \textbf{Sanremo} (LOC)  colocou em  prática medidas para  dar uma nova vida após o processo turbulento e traumático dos últimos anos. O responsável por esse percurso foi \textbf{Amilcare Rambaldi} (PER), um comerciante de flores, que apresentou uma sugestão para dinamizar a povoação: um concurso musical. No início, a administração não estava de acordo com esta ideia, pelo que, só entrou em prática em 1951, depois de ter conhecido \textbf{Angelo Nizza} (PER) e \textbf{Mario Sogliano} (PER), em 1947, que apoiaram esta iniciativa. O então gerente da \textbf{RAI} (ORG) (\textbf{Radio Audizioni Italiane} (MISC)) , \textbf{Giulio Razzi} (PER), permitiu o desenvolvimento de todo o processo  do  que  era  e  é,  hoje  em  dia,  o  \textbf{Festival  de  Sanremo} (MISC).  A  rádio  foi  o \textbf{Media} (LOC) de excelência para a difusão do festival. Nos dias de hoje, esta competição ganhou mais um objetivo: além de elevar a música italiana deu a ver o modo de eleição para o representante italiano no \textbf{FEC} (ORG).

A  primeira  edição  da  \textbf{Eurovisão} (LOC)  deu-se,  como  apresentado  antes,  em  1956,  em \textbf{Lugano} (LOC),  na  \textbf{Suíça} (LOC).  Este  primeiro  concurso  contou  com  a  participação  de  sete  países, considerados países fundadores: \textbf{Alemanha} (LOC), \textbf{Bélgica} (LOC), \textbf{França} (LOC), \textbf{Itália} (LOC), \textbf{Luxemburgo} (LOC), \textbf{Países Baixos} (LOC) e \textbf{Suíça} (LOC). Além da sua participação ativa, contamos com três participantes passivos, isto  é,  que  somente  fizeram  a  transmissão  de  toda  a  competição.  Foram  eles  \textbf{Áustria} (LOC), \textbf{Dinamarca} (LOC) e \textbf{Reino Unido} (LOC). Por ser também a primeira vez em que ocorria, o formato nos anos  seguintes,  foi-se  alterando.  O  formato  inicial  continha  a  apresentação  de  duas canções  por  parte  de  cada  país,  devendo  ser  na  língua  nacional  ou  numa  das  línguas nacionais desse mesmo povo. No final, a \textbf{Suíça} (LOC) saiu vitoriosa com a atuação de \textbf{Lys Assia} (PER), com a canção \textbf{Refrain} (PER) , cantada em francês. A apresentação de duas canções permitia a cada  país  uma  duplicação  das  suas  hipóteses,  relativamente  a  uma  vitória,  visto  que  o formato  inicial  continha  uma  programação  de  uma  hora  e  quarenta  minutos.  Como  o \textbf{Festival  de  Sanremo} (MISC),  foi  através  da  rádio  que  foi  feita  a  primeira  difusão  desta competição.

Para analisarmos as fases de desenvolvimento ou aprimoração do \textbf{FEC} (ORG) devemos refletir sobre as regras que vigoram ou que, por sua vez, já vigoraram, no concurso. Estas regras são, normalmente, aceites por três partes: os organizadores do festival do ano em questão, a \textbf{EBU} (MISC) e o \textbf{Comité Executivo} (ORG) (\textbf{Reference Group} (MISC)) encarregado pela competição.

O formato do concurso, sistema de votação, a apresentação dos resultados das votações às  audiências,  os  valores  do  concurso  e  qual  o  processo  de  divulgação  e  transmissão através dos \textbf{Media} (LOC) (televisão, rádio ou plataformas digitais) são algumas das questões que são definidas anualmente por estas partes. Contudo, alguns pontos continuam iguais, tais como  a  duração  e  a  originalidade  das  canções.  Uma  regra  que  foi  instaurada  em  dois períodos diferentes entre 1966 a 1972 e entre 1978 a 1998, foi a obrigatoriedade de atuar na respetiva língua nacional. Esta regra sofreu, ao longo dos anos, diversas alterações que vieram assinalar uma prevalência do idioma anglo-saxónico, uma vez que foi permitido a inserção de até 25\% de outro idioma, tendo essa percentagem aumentado para 40\% em anos  mais  recentes.  Ao  atuarem  com  uma  canção  na  língua-mãe,  cada  país  tem  a possibilidade de particularizar a sua apresentação. Contudo, devemos ter em conta que alguns  membros  da  \textbf{Eurovisão} (LOC)  têm,  no  seu  país,  mais  do  que  uma  língua  oficial.  Esta regra, dependendo do ponto de vista individual de cada pessoa, pode revelar-se como algo proveitoso pelo facto de cantarem na língua mãe, mas também como algo que lhes poderá retirar alguma atratividade e compreensão pelos outros povos concorrentes. Nos períodos anteriormente  apontados,  percebemos  que  a  diversidade  linguística  tende  a  contribuir para arrecadar bons resultados. O fenómeno de 1969, em que obtivemos quatro primeiros classificados, com apresentações em quatro línguas diferentes, demonstra esse carácter da diversidade linguística e cultural que rege o fenómeno eurovisivo.

Com  o  fim  da  obrigatoriedade  dessa  regra,  começamos  a  observar  o  fenómeno contrário. Em 1999, nota-se um crescimento exponencial do número de atuações que se apresentavam  em  língua  inglesa,  um  fenómeno  que  pode  ser  caracterizado  como  uma monopolização, em certo ponto, das participações. Em 1998, passamos de três atuações em língua inglesa, para doze. Muitas podem ser as razões, uma vez que esta é cada vez mais falada no mundo. O que significa que uma atuação nesse idioma poderia permitir um maior alcance ao nível das apresentações que cada país tem. Há que salientar que, em ambos os anos, o número de participantes total era de 25 e 23, respetivamente.

Apesar  de,  nos  dias  de  hoje,  o  número  de  participações  rondar  os  mesmos números,  houve  uma  crescente  participação  de  países  nos  primeiros  vinte  anos  da competição havendo, como consequência, uma necessidade de aprimorar todo processo de  seleção  e  decisão  no  que  toca  ao  vencedor  do  \textbf{FEC} (MISC).  Consideraremos  fases  de desenvolvimento, quando uma alteração  considerável modifica o formato do concurso.

Como já tivemos oportunidade de referir, o formato de estreia englobava a participação de cada nação com duas canções, situação que ocorreu, única e exclusivamente, em 1956.

A partir de 1957, foi-nos apresentado parte do formato que reconhecemos como atual.

Contudo, as semifinais só foram instauradas anos mais tarde. Estas funcionam como uma pré-seleção de cada participação, isto é, uma fase de qualificação visto que o concurso tinha, cada vez mais, países a participar. O primeiro ano em que encontramos evidência desta ramificação é em 1996, passando a integrar o formato final a partir de 2004. Com a criação  das  semifinais,  o  fenómeno  eurovisivo  passou  a  ocupar  a  atenção  do  público durante  uma  semana  completa,  sendo,  normalmente,  a  1  semifinal  na  terça-feira,  a  2 semifinal na quinta-feira e a final no sábado. Mas, nem todos os países integrantes passam pelas semifinais.

Em  2008,  de  forma  a  recompensar  certos  países  pelos  contributos  financeiros dados, não só ao \textbf{FEC} (ORG), mas também à \textbf{EBU} (MISC), foi criado o grupo dos \textbf{Big 4} (MISC) que englobava a  \textbf{Alemanha} (LOC),  \textbf{Espanha} (LOC),  \textbf{França} (LOC)  e  o  \textbf{Reino  Unido} (LOC).  Com  a  reentrada  da  \textbf{Itália} (LOC),  em  2011, passamos a deter os \textbf{Big 5} (MISC) , que, ainda hoje, integram o formato europeu. O efeito que esta formação detém é garantir um lugar, na final da \textbf{Eurovisão} (LOC), a cada um destes participantes.

Um membro que podemos considerar móvel é o vencedor de cada edição. Graças à sua vitória,  é-lhe  oferecido  um  lugar  na  final  do  concurso  do  ano  seguinte,  ao  lado  dos membros  do \textbf{Big  5} (MISC) .  Caso  ocorra  um  caso  semelhante  ao  do  ano  de  2022,  em  que  o vencedor anterior era um dos membros dos \textbf{Big 5} (MISC) , ficam somente estes cinco apurados para a final, sem necessidade de pré-qualificação em uma das semifinais.

Também  o  sistema  de  votação  foi-se  alterando  durante  cada  ano,  e  em  cada década. Ao longo dos mais de 60 anos de competição, conseguimos encontrar cerca de oito formas de votação. Na primeira edição, dois júris de cada país integrante dariam dois pontos à sua canção favorita. Com a adição de novos participantes, entre 1957 e 1961, passamos a deter dez jurados a pontuar entre 1 e 10 as suas canções favoritas. Até 1975, poucos foram os elementos que sofreram diferença a este nível, sendo 10 o máximo de pontos possíveis.  De 1975 a 1996, cada país  participante teria um  membro jurado que atribuiria pontos às dez melhores canções. A primeira vez que o televoto foi instaurado foi entre 1997 a 2003, através do qual os países poderiam escolher ter membros do júri ou  televoto  para  decidir  como  pontuar  as  melhores  participações.  Entre  2004  e  2008, todos  os  países  utilizaram  o  televoto  como  modo  de  votação.  De  2009  a  2015, conhecemos o início do modelo de votação que é utilizado nos dias de hoje, havendo uma combinação  50/50,  tanto  do  júri,  como  do  televoto.  A  partir  de  2016,  é  possível reconhecer o modelo atual, composto pelos elementos anteriores, mas com dois sets de votação, sendo os votos dos jurados anunciados antes do televoto. Podemos realçar que os votos do júri são dados por um representante de cada nação, após a deliberação de cada grupo de votantes, escolhidos pelo próprio país, ligados de algum modo à cultura. Caso existam  irregularidades  ou  falhas  na  comunicação,  é  o  painel  principal  do  \textbf{FEC} (MISC)  que preside  à  final  da  competição.  Em  caso  de  empate  a  nível  de  pontos,  deverão  ser averiguados quais os países que pontuaram os respetivos países empatados. Se mesmo assim,  ainda  persistir  um  empate,  o  segundo  desempate  deverá  relacionar-se  com  os países que pontuaram 12 pontos aos países empatados, sendo em contagem decrescente até o desempate final.

Apesar  de  pouco  comum,  é  possível  um  país  não  pontuar.  Na  história  da \textbf{Eurovisão} (LOC), esta situação ocorreu em 1962, quando quatro países não receberam qualquer voto. Podemos apontar também 1967, 1970,1998, 2003, 2015 e 2021, entre outros anos.

Consoante as diferentes amizades ou inimizades entre os diferentes países, as pontuações dos jurados podem seguir uma linha de pensamento, isto é, de pontuar não só porque a música  se  apresenta  como  interessante,  mas,  de  igual  modo,  porque  existe  um  elo  de empatia  e  de  entretenimento  entre  esses  países.  Esta  situação  também  pode  ocorrer quando pensamos em tradições ou elementos culturais que liguem diferentes povos, como a  língua,  as  melodias,  entre  outros  aspetos.  Estes  acontecimentos  podem  levar  a  uma especulação a nível da pontuação que cada país atribui. Na edição mais recente do \textbf{FEC} (ORG), em  2023,  o  sistema  de  votação  sofreu  algumas  alterações,  tendo  \textbf{Martin  Österdahl} (PER), \textbf{Supervisor Executivo do FEC} (MISC), anunciado que: \textbf{In} (MISC) 2023 only \textbf{Eurovision Song Contest} (MISC) viewers will decide \textbf{which countries make it to the Grand Final and, reflecting the global impact of the event} (ORG), \textbf{everyone watching the} (MISC) show, \textbf{wherever they live in the world} (MISC), \textbf{can cast their} (MISC) votes for their favourite songs (\textbf{Eurovision} (MISC), 22 novembro de 2022: s/p).

Não só mas também os escolhidos em ambas as semifinais resultaram da votação dos  telespectadores,  algo  que  já  não  fazia  parte  da  realidade  eurovisiva  devido  à combinação entre a votação pública e a do júri. Contudo, a grande inovação apresentada relaciona-se com a votação online que veio permitir um alcance ainda maior da votação, podendo  diversificar  a  escolha  e,  inclusive,  apresentar-se  distinta  das  prospeções inicialmente apresentadas. Os emigrantes encontram, assim, um meio que lhes possibilita apoiar a sua terra natal e os demais concorrentes. Somente na final contamos com os votos do júri, que serão dados na noite anterior à realização da final.

Desde  o  seu  início,  a  \textbf{Eurovisão} (LOC)  funcionou  como  um  canal  de  partilhas  no  que concerne  à  cultura.  Um  elemento  que  se  mostra  relevante,  quando  se  apresenta  um determinado país, são os conhecidos postcards , ou postais, que permitem ao telespectador um maior conhecimento sobre o artista e o país que este irá representar. Nesse sentido, são apresentadas imagens de lugares emblemáticos do país em que decorre o concurso nesse ano assim como do local de origem do intérprete. Em 2023, coube ao \textbf{Reino Unido} (LOC) apresentar  o  \textbf{FEC} (MISC)  e,  devido  às  repercussões  da  \textbf{Guerra} (MISC)  na  \textbf{Ucrânia} (LOC),  encontramos  três representações distintas.

É de salientar que, nestes mesmos postcards , encontramos sítios emblemáticos  de  cada  país  em  que  se  centra  o  \textbf{FEC} (MISC).  Sendo  assim,  a  primeira representação está ligada ao país vencedor da edição anterior, a \textbf{Ucrânia} (LOC), mostrando locais https://www.youtube.com/watch?v=mpAhfmrBxXc - Composição dos postcards da edição de 2023 do \textbf{FEC} (ORG).

como  o  \textbf{Arboretum} (MISC)  \textbf{Sofiyivka} (PER)  em  \textbf{Uman} (LOC),  \textbf{Lviv  Town  Hall} (LOC)  em  \textbf{Leviv} (LOC)  ou  \textbf{St} (MISC).  \textbf{Sophia's Cathedral} (PER) em \textbf{Kiev} (LOC). A segunda representação leva-nos a locais do país que se encontra a realizar o \textbf{FEC} (ORG) nesse mesmo ano, onde foi possível visualizar locais como \textbf{Stonehenge} (LOC) em  \textbf{Salisbury} (LOC),  a  \textbf{Ely  Cathedral} (ORG)  em  \textbf{Cambridgeshire} (LOC)  ou  \textbf{Liverpool  Central  Library} (LOC)  em \textbf{Liverpool} (LOC).  A  terceira  representação  centra-se  no  país  em  questão.  Nesta  última  edição foram  mais  de  100  os  locais  utilizados  para  a  criação  deste  elemento  identificativo  do festival. De um certo ponto de vista, era dado um palco a cada povo que, em condições justas  e  igualitárias,  podiam  mostrar  o  que  de  melhor  e  icónico  os  representa.  Como salientado pela \textbf{EBU} (MISC), no seu portal da \textbf{internet} (MISC) , onde explica todo o desenvolvimento do \textbf{Festival} (LOC),  o  contínuo  crescimento  proporcionado  à  \textbf{Eurovisão} (LOC)  permitiu  que  não  só  a \textbf{Europa} (LOC), mas também o mundo, pudessem estar ligados a esta rede de contactos culturais que nos são apresentados através da \textbf{"} (ORG)televisão nacional\textbf{"} (MISC). Este mesmo fenómeno pode permitir que olhemos, inclusive, para o \textbf{FEC} (ORG) como se de um franchise se tratasse. O termo franchise , originário do francês, mas conhecido pelo seu valor enquanto estrangeirismo à língua inglesa, pode definir-se como "caractère de ce qui est franc […] droit que limitait l'autorité  souveraine  au  profit  \textbf{d'une  ville} (MISC),  d'un  corps,  d'un  \textbf{individu} (LOC)"  (\textbf{Robert} (PER),  1983: 466) . Por se ligar ao termo franquia, este ganha aspetos financeiros e legais. Ao longo dos mais de 60 anos em que esta competição vigora, conseguimos enumerar diferentes manifestações  de  interesse,  por  distintas  partes  do  globo,  em  ter  o  "seu"  \textbf{Festival Eurovisão  da  Canção} (MISC).  Este  interesse  pode  ter  fundamentos  diferenciados  entre  si, contudo, todos partem do mesmo ponto de base, isto é, celebrar a identidade cultural de uma nação. A \textbf{Eurovisão} (LOC) é uma marca característica do continente europeu, sendo o seu começo  um  marco  histórico,  devido  às  situações  apontadas  anteriormente.  Pela  sua dimensão transcontinental, o formato em que poderíamos encontrar uma inspiração ou, por exemplo, um franchise , iria necessitar de ajustes. Em 2022, deu-se o início desta ideia: \textbf{The Eurovision Song Contest is the biggest live music event in the world} (MISC), and \textbf{it is a great honour  to  bring} (MISC)  \textbf{it  to  the  United  States  and  transform  our  dream  of American  Song Contest} (MISC) into a reality on \textbf{NBC} (ORG). \textbf{It's awe-inspiring to stand at the beginning of an American} (MISC) legacy and \textbf{it doesn't} (MISC) get more exciting than that (\textbf{Eurovision} (MISC), s/d: s/p).

\textbf{https://www.ebu.ch/about/history} (MISC) - \textbf{Link} (MISC) para o portal da internet da \textbf{EBU} (MISC).

O \textbf{American  Song  Contest} (MISC) ,  como  observamos  pelo  nome,  é  feito  nos  \textbf{Estados Unidos da América} (LOC). Pretende captar o foco da \textbf{Eurovisão} (LOC) e adaptá-lo à sua realidade, isto é,  a  cada  um  dos  estados  que  os  integram.  Não  só  o  seu  alcance  foi  adaptado,  como também o seu formato. Enquanto no formato europeu contamos com três grandes dias, ambas as semifinais e a final, o formato americano apresentava uma série de episódios, como se de um programa dito comum se tratasse. A parceria entre os produtores "da casa" europeia  e  ambos  o \textbf{Propagate} (LOC) e  a \textbf{Universal  Television  Alternative  Studio} (ORG) mostrou-se fulcral para o bom resultado que teve o \textbf{Festival Americano da Canção} (ORG). Muitos foram os países  que  transmitiram  este  concurso,  um  gesto  que  promove  a  igualdade  de oportunidades,  no  que  concerne  as  demonstrações  televisivas  culturais,  como  por exemplo,  \textbf{Portugal} (LOC),  \textbf{Austrália} (LOC)  e  \textbf{Alemanha} (LOC).  Contudo,  com  base  em  informações  mais recentes, na edição do \textbf{FEC} (ORG) em 2023, \textbf{Martin Österdahl} (PER) confirmou que \textbf{"The American Song  Contest} (MISC)  was  cancelled  after  its  first  season  on  \textbf{NBC} (ORG)  due  to disappointing ratings and lack  of  impact  on  the  music  charts"  (\textbf{Eurovisionfun} (LOC),  2023:  s/p),  estando  os mesmos  à  procura  de  novos  parceiros  para  a  realização  de  um  evento  inspirado  no fenómeno eurovisivo. A expansão do formato eurovisivo às regiões ou países que não o compõem  permite  elevar  o  potencial  do  concurso  enquanto  manifestação  cultural.  Em resultado disso, e após a edição de 2022 do \textbf{FEC} (ORG), dois novos formatos foram anunciados igualmente. Em 2023, irá ocorrer pela primeira vez o \textbf{Eurovision Song Contest Canada} (ORG) que  promete  encontrar  o  melhor  representante  de  entre  10  províncias  canadianas.  Este novo formato poderá vir a seguir o modelo americano, visto que se restringe a um único país. Também o \textbf{Eurovision Song Contest Latin America} (ORG) integrará a família da \textbf{Eurovisão} (LOC).

Possivelmente, um formato mais semelhante ao europeu que pretende, \textbf{Eurovision  Song  Contest  Latin  America} (MISC)  marks  the  latest  development  in  the  global expansion of the Eurovision Song Contest as it simultaneously elevates \textbf{the legacy of the brand while maintaining local cultural representation} (MISC) (\textbf{Eurovision} (MISC), 2022: s/p).

Isto é, depreende-se que a \textbf{Eurovisão} (LOC) funciona como um palco cultural, importante para  os  diferentes  países  que  integram  este  concurso,  pelo  que,  podem  estar  mais próximos de um público, presencial ou digitalmente, que pode vir a interessar-se pelas suas tradições ou modos culturais.

Se quisermos seguir esta ideia de franchise , podemos também pensar num evento que, apesar de constituído por países europeus, não parece conter, em certas nações, até à data tanta relevância, sendo ele o \textbf{Festival Eurovisão da Canção Júnior} (MISC) (\textbf{FECJ} (ORG)). Com início  em  2003,  o \textbf{Junior  Eurovision  Song  Contest} (MISC) ,  em  \textbf{Copenhaga} (LOC),  na  \textbf{Dinamarca} (LOC), permitiu que as crianças pudessem fazer parte de um evento semelhante ao da \textbf{Eurovisão} (LOC).

Além de fomentar o gosto pela música e pela cultura, são os vencedores, muitas vezes, convidados a participar no \textbf{"Grande Prémio} (MISC)", ou seja, no \textbf{FEC} (MISC), quando tiverem mais idade.

Podemos  salientar  o  nome  de  \textbf{Irina  Khechanovi} (LOC),  representante  da  \textbf{Geórgia} (LOC)  na  edição número 67 do \textbf{FEC} (ORG), que já havia participado no \textbf{FECJ} (ORG). Procuraremos ainda elencar quais os pontos em comum entre a versão original e a versão para jovens.

O  \textbf{FECJ} (MISC),  ao  seguir  os  ideais  e  propósitos  do  \textbf{FEC} (MISC),  coloca  ao  serviço  dos  mais jovens uma plataforma de diversão e representação de cada nação. Apesar de com menos edições do que o original, ao longo dos anos tem vindo a obter mais destaque. É claro que, apesar de ser uma extensão do \textbf{FEC} (ORG), nem todos os aspetos são semelhantes entre si.

Ao  contrário  da  regra  da  língua,  que  passou  por  diversas  fases  no  \textbf{FEC} (MISC),  no  \textbf{FECJ} (MISC),  os participantes  são  obrigados  a  apresentar  uma  música  que  seja  escrita  e  cantada  na  sua língua-mãe, podendo contar com a ajuda de adultos, desde 2008, na composição musical.

A própria data de realização do \textbf{FECJ} (ORG) pretende manter vivo o espírito eurovisivo ao longo do ano, visto que o \textbf{FEC} (ORG) se realiza em maio e o \textbf{FECJ} (ORG), no final do ano, regularmente em dezembro. Parte  dos jovens que competem  neste concurso encontram  os seus pares na versão original do \textbf{Festival} (PER). Porém, existe um país que se apresenta distinto, sendo ele o \textbf{Cazaquistão} (LOC), que conta com a sua entrada no \textbf{FECJ} (ORG). Na sua história participativa, apesar de ser um país da periferia europeia, conta com o segundo lugar em dois anos, após ter entrado pela primeira vez em 2018.

Não  só  mas  também  em  solo  europeu  e  além-fronteiras,  é  possível  denotar inspirações  deste  formato  nos  mais  variados  campos.  No  \textbf{Velho  Continente} (LOC),  podemos referir  a  nível  mais  local  o  \textbf{Bundesvision} (LOC),  uma  adaptação  da  \textbf{Eurovisão} (LOC)  aos  estados alemães, a nível mais geral o \textbf{Eurovision Young Dancers} (ORG), \textbf{Eurovision Young Musicians} (MISC), \textbf{Eurovision  Dance  Contest} (MISC)  e  \textbf{Magic  Circus  Show} (MISC).  Atualmente,  somente  o  \textbf{Eurovision Young Musicians} (MISC) ainda se realiza. A nível extra-europeu, conseguimos encontrar no lado asiático, o \textbf{Turkvision Song Contest} (ORG), \textbf{Bala Turkvision Song Contest} (ORG) e os diferentes \textbf{Asia- Pacific Broadcasting Union} (ORG) (\textbf{ABU} (LOC)) \textbf{Song Festivals} (PER), \textbf{ABU Radio Song Festival} (LOC), \textbf{ABU TV Song Festival} (LOC), \textbf{ABU International Dance Festival} (LOC) e \textbf{ABU Song Contest} (LOC). Estas recriações e inspirações promovem a ideia de franchise e levam o core do \textbf{FEC} (ORG) para além do seu normal raio de influência.

Atualmente, o \textbf{FEC} (ORG) busca reencontrar o seu papel de atratividade junto dos mais novos e das comunidades digitais. É através destas mesmas plataformas que parece existir um maior fator de disparidade quando pensamos nas votações do público versus a votação dos jurados de cada nação. \textbf{Plataformas} (PER) que apresentam um algoritmo mutável, como por exemplo  o tiktok ,  permitem  a  constante  atualização  daquilo  que  é  apresentado  ao consumidor, podendo criar novos padrões musicais ou expetativas para o público, tema que abordaremos num outro ponto, quando reflectirmos sobre algumas atuações na edição número  sessenta  e  seis  do  \textbf{FEC} (MISC).  \textbf{Analisaremos  de  que  forma  estes  fatores  poderão eventualmente alterar} (MISC) uma votação que seria vista como provável.

2.2.

Canção a canção: intervenientes, língua e vestuário Nos  sessenta  e  sete  anos  de  história  eurovisiva  até  à  data,  é  possível  encontrar diversos aspetos que permitem distinguir os intervenientes em palco, seja o número de intervenientes,  a  sua  língua  ou  outros  meios  utilizados.  Neste  sentido,  e  por  existirem inúmeras edições, optámos, na nossa investigação, por fazer uma abordagem comparativa entre a primeira edição do \textbf{FEC} (MISC), realizada em 1956, e a sua edição mais recente, em 2023, refletindo os processos e recursos envolvidos ao longo destes anos.

Como  apresentado  anteriormente,  em  1956,  participaram  pela  primeira  vez  no \textbf{FEC} (ORG) sete países (\textbf{Bélgica} (LOC), \textbf{França} (LOC), \textbf{Alemanha} (LOC), \textbf{Itália} (LOC), \textbf{Luxemburgo} (LOC), \textbf{Países Baixos} (LOC) e \textbf{Suíça} (LOC)), sendo esses considerados os membros fundadores desta competição. Apesar de atuarem duas vezes, com músicas distintas, constatamos que a maioria dos participantes recorreu à mesma língua em ambas as demonstrações, exceto a vencedora, oriunda da \textbf{França} (LOC), \textbf{Lys Assia} (PER), que cantou nas línguas francesa e alemã, duas das línguas oficiais suíças. Contudo, nessa primeira edição, a língua francesa foi o idioma mais utilizado pelos participantes, incluindo  \textbf{França} (LOC),  \textbf{Luxemburgo} (LOC)  e  \textbf{Bélgica} (LOC).  Ao  contrário  dos  demais,  os  mesmos representantes suíços e \textbf{luxemburgueses} (LOC) executaram as duas performances, o que poderia ajudar o público a uma identificação dos intérpretes do país de origem de uma maneira mais  fácil.  Além  dos  participantes  que  disponibilizaram  logo  através  dos  meios radiofónicos, as suas atuações, três outros países fizeram o mesmo, entrando oficialmente no ano seguinte na competição. Referimo-nos a  \textbf{Áustria} (LOC), \textbf{Dinamarca} (LOC) e \textbf{Reino  Unido} (LOC). O concurso  decorreu  durante  um  único  dia,  demorando  igualmente  menos  tempo,  como explicado no capítulo anterior.

Refletindo agora sobre o século \textbf{XXI} (MISC) e a última edição do \textbf{FEC} (MISC), em 2023, face ao formato  que  havia  sindo  apresentado  na  data  de  estreia  deste  concurso,  é  possível encontrarmos inúmeras alterações. Graças a diversas dinâmicas históricas ocorridas em solo europeu, notamos um incremento no número de países participantes: de 7 países em 1956,  assistimos,  em  2023,  a  37  países  em  concurso.  Apesar  de  não  contar  enquanto participante, foi facultada a informação de que, passado trinta anos, o \textbf{Luxemburgo} (LOC), país fundador do \textbf{FEC} (ORG) e vencedor do concurso por 5 vezes até à data, voltará a participar na edição de 2024 do festival. Em virtude das regras impostas no que concerne às canções apresentadas,  percebemos  que  vinte  e  três  das  trinta  e  sete  apresentaram  músicas  em língua inglesa. Apesar de algumas canções conterem passagens em inglês, eram cantadas em outro idioma característico do país  que estaria a ser representado. Porém,  \textbf{Albânia} (LOC), \textbf{Croácia} (LOC),  \textbf{Eslovénia} (LOC),  \textbf{Espanha} (LOC),  \textbf{Finlândia} (LOC),  \textbf{França} (LOC),  \textbf{Itália} (LOC),  \textbf{Moldávia} (LOC)  e  \textbf{Portugal} (LOC) apresentaram-se  em  palco  com  composições  musicais  na  sua  língua-mãe.  Destes  nove países,  se  pretendermos  analisar  de  que  forma  se  saíram  a  nível  de  pontuação  em correlação  com  a  utilização  do  idioma  identificativo,  percebemos  que  somente  dois, \textbf{Finlândia} (LOC) e \textbf{Itália} (LOC), arrecadaram um lugar no top 10. A canção vencedora deste ano foi a da \textbf{Suécia} (LOC), com atuação de \textbf{Loreen} (LOC), que já havia representado o seu país em 2012, tendo sido consagrada vencedora nesse mesmo ano.

\textbf{Ligadas} (PER) à cultura do país, existem participações que os fãs eurovisivos consideram icónicas e que, mesmo não tendo sido vencedoras, fazem parte da história do festival. A título de exemplo, podemos referir um dos nomes mais populares desta competição, no século XX, logo dois anos após o início do \textbf{FEC: Domenico Modugno} (MISC), em representação da \textbf{Itália} (LOC), apesar de ter arrecadado o terceiro lugar. Nesse mesmo ano, compôs uma das músicas eurovisivas mais cantadas. Numa época em que a rádio era considerada o meio de  transmissão  por  excelência,  o  ritmo  animado  conquistou  os  fãs.  No  ano  seguinte, \textbf{Modugno} (PER) viria a ganhar o \textbf{Grammy} (MISC) para música do ano e para gravação do ano. A música \textbf{Nel  Blu  Dipinto  Di  Blu} (MISC) é  a  única  composição  eurovisiva  a  arrecadar  estes  mesmos prémios,  demonstrando  que  o  festival  poderia  e  pode  funcionar  como  uma  rampa  de lançamento para os artistas. Se recuarmos a 1974, recordamos os \textbf{ABBA} (ORG), a banda sueca que conquistou o coração dos fãs e que se tornou numa das bandas mais acarinhadas dos anos  70,  a  nível  mundial.  É  de  sublinhar  que  no  século  \textbf{XXI} (LOC)  continuam  a  provocar  a atenção do público, tendo em conta os vários filmes e documentários que se continua a produzir sobre essa banda. Aquando da sua atuação no \textbf{FEC} (MISC), o vestuário utilizado colocou em palco um estilo descontraído, brilhante e livre que viria a caracterizar toda a década: \textbf{Everything} (LOC) changed as \textbf{Agnetha and Frida} (ORG) bounced down the stage in their gloriously of- the-moment synthetic outfits and glittery make-up. (…)  \textbf{Reworking the song's Swedish lyrics} (MISC) into \textbf{English} (PER) was another master-stroke, \textbf{broadening the glam-pop stomper's appeal far further than if it had remained in its native language} (MISC). \textbf{None of that would likely have mattered} (MISC),  though,  if  "\textbf{Waterloo} (MISC)"  hadn't  been  such  a  strong  song,  and  the  \textbf{ABBA Eurovision} (MISC) performance so assured. (…) \textbf{It remains one of the classic Eurovision records} (MISC), \textbf{illustrated by its subsequent hike up the charts} (ORG) (\textbf{Elliot} (LOC), 2023: s/p).

Já neste século e talvez pelo constante despoletar do recurso aos meios digitais, reconhecemos  determinados  artistas  vencedores  desta  competição.  \textbf{Nomes} (PER)  como  os  de \textbf{Alexander Rybak} (PER), representante da \textbf{Noruega} (LOC) em 2009, \textbf{Loreen} (LOC), representante da \textbf{Suécia} (LOC) em 2012 e 2023, \textbf{Emmelie de Forest} (PER), representante da \textbf{Dinamarca} (LOC) em 2013, e \textbf{Conchita Wurst} (PER),  representante  da  \textbf{Aústria} (LOC)  em  2014,  permanecem  na  memória  eurovisiva  mais recente pelo seu papel de animação ou pelos ideais que defendiam. Mas, como vimos no caso de \textbf{Domenico Modugno} (PER), não só de vencedores se escreve a história deste festival.

Apesar de muitos participantes não terem sido consagrados vencedores são, hoje em dia, relembrados. Nos anos mais recentes do \textbf{FEC} (ORG), \textbf{Mahmood} (MISC), da \textbf{Itália} (LOC), \textbf{Go\_A} (LOC), da \textbf{Ucrânia} (LOC) e \textbf{Konstrakta} (LOC), da \textbf{Sérvia} (LOC), são alguns dos nomes que, nos seus respetivos anos, conseguiram um papel de diferenciação e de distinção entre os seus pares.

Da  mesma  forma,  como  encontramos  novos  artistas  que  buscam  lançar-se  no mundo da música através das portas eurovisivas, alguns são aqueles que, depois da fama, integram este mesmo concurso para manter vivo o seu papel na música, como foi o caso dos  \textbf{The  Rasmus} (MISC),  pela  \textbf{Finlândia} (LOC),  e  os  \textbf{Hooverphonic} (ORG),  pela  \textbf{Bélgica} (LOC).  \textbf{Fazer} (LOC)  parte  de  um evento com a proporção da do \textbf{FEC} (ORG) é uma oportunidade para lançar qualquer artista, como já  tivemos  oportunidade  de  referir.  \textbf{Vencedores} (LOC)  ou  não,  os  mesmos  devem  saber capitalizar a sua participação para que assim possam ganhar ainda mais reconhecimento e fama. Em certas situações, é a música que conquista os fãs, tornando-se numa das mais ouvidas ao longo das diferentes plataformas digitais, como \textbf{Spotify} (ORG), \textbf{Apple Music} (ORG), \textbf{Youtube} (MISC) , entre outras. Tal foi o caso de \textbf{Duncan Laurence} (PER), em representação dos \textbf{Países Baixos} (LOC), que totaliza  mais  de  935000000 streams no  \textbf{Spotify} (LOC),  "securing  its  position  as  \textbf{the  ultimate streaming  champion  among} (MISC)\textbf{Eurovision  winners} (MISC)"  (\textbf{Vranis} (PER),  2023:  s/p).  Outro  exemplo pode ser encontrado com os vencedores da edição de 2021 do \textbf{FEC} (ORG), os \textbf{Måneskin} (LOC), cuja vitória levou a que conquistassem fãs além do território europeu. Podemos considerar que o  feito  eurovisivo  alargou  o  raio  de  influência  desta  banda,  especialmente  através  das redes sociais, levando o sonho da banda, que havia participado no \textbf{X Factor} (MISC) italiano anos antes,  além  fronteiras.  "\textbf{While  Måneskin} (MISC)'s  inimitable  swagger  have  \textbf{led  to  a  recent international meteoric rise} (MISC), \textbf{the Italian} (MISC) quartet have tapped the glitter and grime of rock's glory days" (\textbf{Phillips} (PER), 2023: s/p), conduzindo à sua nomeação, em 2023, para o \textbf{Grammy de  Artista  Revelação} (MISC).  Desde  a  sua  vitória,  a  referida  banda  tem  estado  em tour nos diferentes  continentes,  cativando  aqueles  por  onde  passam.  Numa  sociedade  em  que procuramos  a  igualdade,  dentro  das  nossas  diferenças,  sejam  elas  étnicas,  sociais  ou sexuais,  as  redes  sociais  e  as  plataformas  digitais  funcionam  como  um  meio  de transmissão das ideias e ideais. A fluidez de género entre os quatro membros da banda e a aceitação de todas as preferências individuais destacam-se como elementos que marcam e fundamentam a presença dos \textbf{Måneskin} (PER) na agenda pública. Muitos são os projetos que advieram da vitória dos mesmos no \textbf{FEC} (ORG), provando, uma vez mais, que existe o papel de lançamento e o de alavancar o potencial de um artista que compita neste concurso.

Independentemente  do  reconhecimento  que  possa  advir  desta  carreira,  é  no momento em que se apresentam em palco que podemos encontrar a festa de uma cultura levada  ao  seu  ponto  máximo.  As  opiniões  podem  diferenciar-se  e  os  fãs  eurovisivos podem  não  ser  apologistas  da  vertente  cultural  e  tradicional  de  um  determinado  povo, mas,  se  refletirmos  no  papel  que  estas  podem  ter  em  diferenciar-se  das  demais, percebemos que podem tender para ambas as posições. É neste sentido que, pensando nos aspetos culturais que explorámos no capítulo anterior, consideramos pertinente analisar algumas  intervenções  artísticas  que  nos  parecem  relevantes  apresentar  na  nossa investigação.

"Mesmo do ponto de vista linguístico, basta um pouco de ouvido para reconhecer as  diferenças"  (\textbf{Cabral} (PER),  2001:  73),  sendo  essas  diferenças  que  poderão  influenciar  a votação.  \textbf{França} (LOC),  em  2022,  apresentou-se  com  a  música \textbf{Fulenn} (PER) ,  cantanda  num  dialeto francês, o \textbf{bretão} (LOC), da região da \textbf{Bretanha} (LOC), no noroeste francês. Além da apresentação com um elemento tradicional, como é a língua, toda a performance de \textbf{Alvan} (LOC) \&amp; \textbf{Ahez} (LOC) estava envolvida  por  um  carácter  de  mistério  e  lendário  atribuído  àquela  mesma  região.  Foi através  desse  carácter  de  diferenciação  que  foram  selecionados  na  pré-qualificação francesa, \textbf{C'est} (MISC)  vous  qui  décidez .  Quando  questionados  sobre  a  apresentação  de  uma música  composta  por  uma  linguagem  e perfomance característica  de  uma  localização, afirmaram que: \textbf{It} (MISC)'s true that \textbf{when we speak about minority languages} (MISC), \textbf{we often think in terms of the past} (MISC).

\textbf{But we really wanted to show that this language we are representing is actually alive and healthy} (MISC),  and  we  can  mix  \textbf{it  with  electro  music  and  other  influences} (MISC).  \textbf{The  Breton} (MISC) language - we  want  \textbf{it  to  grow} (MISC).  \textbf{It} (MISC)  really  can  \textbf{fit  every  style  of  music} (MISC).  \textbf{That's  what's  so magical} (MISC) with music, we can explore a \textbf{mix of so many} (ORG) different things and that's what we wanted to show - bringing the language to such a great event \textbf{like the Eurovision Song Contest} (MISC) (\textbf{Eurovision} (ORG), 2022: s/p).

A  título  de  informação,  a  canção  apresentada  não  conseguiu  arrecadar  o  lugar desejado, ficando em vigésimo quarto lugar na final, totalizando dezassete pontos, sendo oito  providenciados  pelo  júri  e  nove  pelo  público.  Recuando  no  tempo,  uma  situação semelhante foi a da \textbf{Turquia} (LOC), em  1975, representada por \textbf{Semiha Yanki} (PER),  que  arrecadou apenas três pontos na sua estreia no \textbf{FEC} (ORG). Não obstante, em outros anos, encontramos o fenómeno  inverso.  Ao  entrar  no  festival  em  1973,  \textbf{Israel} (LOC),  nos  seus  primeiros  anos, apresentou-se, única e exclusivamente com uma canção em hebraico. Contudo, foi só na sua sexta e sétima participação que conseguiu sagrar-se vencedor, mantendo a língua no centro  da  atenção.  Nos  últimos  vinte  anos  do  \textbf{FEC} (MISC),  somente  quatro  vencedores arrecadaram este lugar, homenageando o seu país com uma canção na língua-mãe: \textbf{Sérvia} (LOC), em 2007, \textbf{Portugal} (LOC), em 2017, \textbf{Itália} (LOC), em 2021, e \textbf{Ucrânia} (LOC), em 2022. Este resultado poderá explicar a prevalência e, de certo modo, preferência de um artista em se apresentar com uma composição em língua inglesa.

"Quando  as  fronteiras  se  dissolvem,  as  identificações  também  se  confundem" (\textbf{Brandão} (LOC), 2011: 41) e o vestuário que cada indivíduo utiliza pode encontrar e identificar o seu semelhante, como vimos no anterior capítulo. As vestes típicas ou identificativas de um  determinado povo podem  levar a uma identificação dos mesmos,  apesar de não se apresentarem  na  sua  língua  principal.  Ao  longo  dos  anos,  diferentes  trajes  serviram  e marcaram as participações. Dois dos casos mais recentes podem levar-nos a 2003, quando a \textbf{Turquia} (LOC), com \textbf{Sertab Erener} (PER), se apresentou em palco com vestuário semelhante àqueles que encontraríamos na dança do ventre e sonoridades características deste mesmo estilo que  é  facilmente  encontrado  em  diversas  regiões  da  \textbf{Turquia} (LOC),  sagrando-se  assim  nesse mesmo  ano  vencedora  do  \textbf{FEC} (MISC).  Também  \textbf{Buranovskiye  Babushki} (LOC),  em  2012,  em representação  da  \textbf{Rússia} (LOC),  utilizou  vestes  típicas,  com  tonalidades  semelhantes  ao vermelho e padrões clássicos repetitivos, em homenagem  a um dos trajes mais antigos russos, o poneva , arrecadando o segundo lugar, ficando a cento e treze pontos do vencedor desse mesmo ano. Nos anos 80, também é possível encontrar este tipo de representação cultural  com  a  \textbf{Noruega} (LOC)  que,  apesar  da  individualidade,  ficou  aquém  do top 10  desse mesmo ano.

Os instrumentos e os géneros musicais também contribuem para que um país se distinga,  mesmo  quando  não  estão  de  forma  tão  vincada  na  sua  cultura  ou  nas  suas apresentações. Apesar de menos comum, é possível encontrarmos este mesmo feito, em 1997, com a \textbf{Grécia} (LOC) e a \textbf{Turquia} (LOC). Representadas por \textbf{Marianna Zorba} (PER) e \textbf{Sebnem Paker and Group Etnic} (PER), respetivamente, conseguiram distinguir-se dos outros concorrentes através do uso de instrumentos como o krotalon e o baglama . Embora façam parte de famílias musicais distintas, um de percussão e outro de cordas, estes elementos adicionaram outro nível de diferenciação para além da língua e do vestuário.

Em  anos  mais  recentes,  e  com  a  retirada  da  orquestra  ao  vivo  que  poderia acompanhar os diferentes artistas nas suas atuações, em 2004, passou a ser mais difícil encontrar instrumentos que permitissem este efeito distintivo. O caso mais recente deu- se em 2021, com os \textbf{Go\_A} (MISC), banda que já referimos anteriormente, utilizou um sopilka , um  instrumento  de  sopro,  fabricado  em  madeira,  que  nos  transporta  para  locais  na natureza, tendo já sido utilizado em apresentações anteriores da \textbf{Ucrânia} (LOC).

Ainda  que  não  necessitemos  de  combinar  todos  estes  elementos  para  que tenhamos uma representação mais genuína de uma determinada cultura ou de um povo, quando encontramos  a conjugação de todos estes pontos,  nota-se a presença do aspeto cultural,  independentemente  da  sua  classificação.  É  certo  que,  quanto  mais  perto estejamos  de  uma  cultura,  mais  fácil  será  para  o  indivíduo  reconhecer  estes  mesmos componentes.  Um  facto  que  também  pode  surgir  devido  aos  conceitos  previamente associados a cada comunidade, podendo fazer parte de um pré-conceito, em certa medida.

Após  uma  pesquisa  de  entre  os  mais  variados  elementos  anteriormente  listados,  a apresentação de \textbf{Blanca Paloma} (PER), em representação da \textbf{Espanha} (LOC), em 2023, poderia somar todos estes itens. A composição \textbf{EAEA} (ORG), cantada exclusivamente em castelhano, transmite a  energia  associada  ao  flamenco,  algo  característico  um  pouco  por  todo  o  território espanhol.

" \textbf{which is in keeping with Blanca's trademark style of mixing the traditional with the avant-garde} (MISC) and modern pop" (\textbf{Eurovision} (MISC), s/d: s/p). A utilização da cor vermelha e  da  luz  transporta  o  espetador  aos  trajes  típicos  em  que  esta  tonalidade  representa  a paixão e a fugacidade do momento: \textbf{This} (MISC) songstress has a playlist \textbf{full of Spanish} (MISC) influences, as \textbf{well} (ORG) as a \textbf{few} (ORG) songs from across the  globe.  \textbf{Blanca  loves  the  vintage} (LOC)\textbf{Spanish} (MISC)  vibes  of \textbf{La} (MISC)Bomba  \textbf{Gitana} (LOC) by  \textbf{Lola  Flores and Quien Maneja} (LOC) Mi Barca by Remedios Amaya (\textbf{Eurovision} (LOC), 2023: s/p).

Esta  é  uma  das  mais  recentes  manifestações  que  podem  ser  transversalmente associadas  a  um  povo  e  a  uma  cultura  pela  sequencialidade  dos  elementos.  Podemos referir participações anteriores, como a da \textbf{Irlanda} (LOC), em 2007, com \textbf{Dervish} (PER), uma banda de música folk irlandesa que se apresentou em palco utilizando trajes inspirados nas vestes dos  povos  celtas  que  estiveram  presentes  no  que  é  atualmente  considerado  território irlandês. Os trajes associados à letra e à melodia  conduzem o espectador à \textbf{Irlanda} (LOC) dos séculos anteriores, sobretudo à \textbf{Irlanda} (LOC) do século XV e da animação que era apresentada nas cortes reais.

2.3.

\textbf{Tendências} (LOC), manifestações culturais e desafios Embora  o  \textbf{FEC} (MISC)  seja  um  concurso  onde  os  seus  intervenientes  procuram  uma vitória, é normal que encontremos elementos recorrentes e pontos de diferenciação face às demais competições. Como tentámos demonstrar anteriormente, com a não obrigação da  utilização  da  língua  mãe  numa  determinada  apresentação  é  percetível  uma monopolização da língua inglesa. "\textbf{According to Ethnologue} (MISC), \textbf{English is the most-spoken language in the world including native and non-native speakers} (MISC)" (\textbf{Berlitz} (PER), 2023, s/p). Uma afirmação que revela como uma língua pode tornar mais simplificada a compreensão e captação de mais ouvintes e, por conseguinte, tende a arrecadar melhores resultados. Nos últimos  vinte  anos,  dezasseis  dos  vencedores  apresentaram-se  em  palco  com  uma composição  em  língua  inglesa,  dois  dos  quais,  detinham  uma  combinação  entre  esta língua e o idioma tradicional do país que representavam, neste caso, a \textbf{Ucrânia} (LOC). Passámos de  dezasseis  em  vinte  e  seis  participantes,  em  2003,  para  vinte  e  três  em  trinta  e  sete participantes,  em  2023,  que  utilizam  a  língua  inglesa  como  meio  de  expressão.  Esta tendência, que tem vindo a tornar-se mais acentuada com o decorrer do tempo, continuará certamente a revelar-se ainda mais evidente, a não ser que o \textbf{Comité Executivo} (ORG) do \textbf{FEC} (ORG) volte a introduzir a regra da obrigatoriedade de utilização da língua oficial do país que representam.

Parece-nos  pertinente  para  a  nossa  dissertação  refletir  sobre  a  temática  das canções dos vencedores do \textbf{FEC} (ORG). Por já terem existido inúmeras edições, sessenta e sete, centraremos a nossa atenção, a título de exemplo, nos últimos dez anos de competição. A melodia, composição ou até estilo musical pode diferenciar-se dos demais, mas, é no core da  letra  dessas  canções  que  encontramos  uma  possível  razão  que  pode  ter  contribuído para  a  vitória  desse  mesmo  artista.  Mais  de  metade  dos  dez  vencedores  anteriores apresentaram composições que se centram em questões sociais, de insurgência social em certa medida, observando o indivíduo de um ponto de vista exterior, podendo oferecer- lhe motivações e razões que o sustentem. Por exemplo, a composição \textbf{Heroes} (MISC) , apresentada por  \textbf{Mäns  Zelmerlöw} (MISC), em  2015,  relata  como  o  indivíduo  está  rodeado  de  diversas situações que podem dificultar a forma como o mesmo  as encara, nunca esquecendo o poder que este conta dentro de si, afirmando: \textbf{Don't tell the gods I} (MISC) left a mess I can't undo what has been done \textbf{Let's run for cover What if I'm the only hero left?} (MISC)You better fire off your gun Once and forever \textbf{He said go} (MISC) dry \textbf{your eyes} (MISC) And live your life like there is no tomorrow son \textbf{To go sing it like a hummingbird} (MISC) \textbf{The greatest anthem ever heard} (MISC) \textbf{We are the heroes of our time} (MISC) \textbf{But we're dancing with the demons in our minds} (MISC) Também \textbf{Zitti} (PER) e \textbf{Buoni} (PER) , de \textbf{Mäneskin} (PER), traduzida para português como quieto e bom, no sentido de calado e bem-comportado, reflete sobre como não nos devemos conformar com o que é a norma da sociedade, encontrando nos outros semelhanças através das suas diferenças, como abordamos nesta investigação quando refletimos sobre multiculturalidade.

Outro  tópico  que  poderíamos  encontrar  nestes  últimos  vencedores  está relacionado com a temática do amor, podendo este ter sido correspondido ou não. Não só, mas também, as dificuldades que pode advir de uma relação amorosa. Em 2018, \textbf{Netta} (LOC) apresentou-se em palco com a composição \textbf{Toy} (MISC) , uma música que se tornou popular apesar de muitos  questionarem  o que fez  levar  a este fenómeno. A letra afirma-se como  uma elevação de um dos lados da relação amorosa após esta não ter resultado, dizendo \textbf{"I} (PER) am not your toy". No ano seguinte, \textbf{Duncan Laurence} (PER) também se sagrou vencedor com uma canção com uma letra/temática semelhante, ao referir: A \textbf{broken heart is all that's left I'm still fixing all the cracks} (MISC) \textbf{Lost} (MISC) a \textbf{couple of pieces when I carried it} (ORG), \textbf{carried it} (MISC), \textbf{carried it home I'm afraid of all I am My mind feels like} (MISC) a foreign land Silence ringing inside my head Please, carry me, carry me, \textbf{carry me home I've spent all of the love I saved} (MISC) We were always a losing game \textbf{Small-town boy in} (MISC) a big arcade I got addicted to a losing game Se centrarmos a nossa atenção apenas nas composições vencedoras observamos que as temáticas, tais como a do amor, do malogro sentimental e da irreverência perante medidas sociais, são as mais populares entre as mesmas. Contudo, os tópicos das demais canções  tendem  a  encontrar  pontos  em  comum  ou  semelhantes  com  as  vencedoras.  A letra  carrega  um  valor  importante  para  cada  apresentação  singular,  mas,  é  também  a melodia  que  possibilita  que  uma  composição  arrecade  um  bom  lugar  na  tabela classificativa.  Das  canções  que  apresentámos  acima,  sobretudo  as  que  se  encontram ligadas  ao  domínio  sentimental,  apresentam-se  com  uma  melodia  semelhante  a  uma balada, enquanto as demais, ligadas à tomada de posição perante questões sociais exibem uma melodia mais áspera e dura, acompanhando a música das mesmas. Disso pode ser exemplo  a utilização do  estilo  pop rock pelos  \textbf{Mäneskin} (PER),  utilizando o core deste estilo musical  que  é  a  sua  irreverência  e  a  contestação  da  normalização.  No  anexo  5, encontramos as letras completas das canções vencedoras dos últimos dez anos.

A confraternização na \textbf{Green Room} (ORG) e as semelhanças entre os países podem levar a  que  existam  elos  entre  os  mesmos,  promovendo  uma  facilitação  ou  favorecimento  a determinadas nações no voto que evidenciamos como, por exemplo, o voto por amizade.

Este tipo de voto, reconhecido na votação dos jurados, tende a seguir uma tendência. Se um  determinado  país  vota  com  uma  respetiva  pontuação  noutro,  esse  voto  pode  ser repercutido  no  sentido  inverso.  Exemplo  dessa  tendência  poderia  ser  a  ligação  entre \textbf{Chipre} (LOC)  e  \textbf{Grécia} (LOC).  Entre  2016  e  2021,  nas  cinco  edições  do  \textbf{FEC} (MISC),  a  \textbf{Grécia} (LOC)  pontuou  a apresentação  do  \textbf{Chipre} (LOC)  com  o  máximo,  ou  seja,  com  doze  pontos.  Em  2022,  essa pontuação passou para dez pontos. No reverso, entre 2016 e 2022, e apesar de não se ter qualificado na edição de 2018 para a final, \textbf{Chipre} (LOC) pontuou, em todos esses anos, \textbf{a Grécia} (LOC) com doze pontos. Se ficássemos por aqui, poderíamos afirmar que se manteria este voto por amizade. Contudo, no mais recente registo do \textbf{FEC} (ORG), encontramos uma alteração a esta situação. \textbf{A Grécia} (LOC) falhou a sua passagem à final e com a alteração das votações apenas nos podemos referir à final.  Neste sentido, é facto que \textbf{a Grécia} (LOC) pontuou  com somente quatro  pontos  a  participação  de  \textbf{Chipre} (LOC),  indo  contra  aquilo  que  pudemos  observar  nos últimos anos.

Desde  a  sua  fundação,  o  \textbf{Festival  da  Eurovisão} (MISC)  tem-se  apresentado  como  um movimento  apolítico  para  que  funcionasse  como  um  território  neutro  no  \textbf{pós-Guerras Mundiais} (MISC). Porém, por ser composto por diferentes gentes e resultado de uma comunidade com interesses distintos, a sua imparcialidade aos confrontos internacionais pode ter-se esbatido ao longo dos mais de sessenta anos de competição: \textbf{The main aim of cultural diplomacy is to strengthen diplomatic relations and improve the international image of a country in order to advance diplomatic} (MISC), political, and economic ties and interests, as \textbf{well} (ORG) as to maintain and normalize bilateral ties in time of tension (\textbf{Borić} (LOC), 2017: 226).

Contudo, desde 1958, conseguimos apontar certos aspetos que tendem a provar a existência de uma discordância ou um posicionamento, em resultado de políticas ou ações pouco  amigáveis,  não  só  entre  os  participantes,  mas  também  entre  outros  países periféricos  ao  \textbf{Velho  Continente} (LOC).  O  caso  mais  recente,  que  veio  a  culminar  em  duas edições repletas de suposições, foram as de 2022 e 2023. Como abordámos anteriormente, a invasão da \textbf{Ucrânia} (LOC)  aconteceu  meses  antes do  início do festival  da \textbf{Eurovisão} (LOC). Nesse momento,  os  dois  atores  principais,  \textbf{Ucrânia} (LOC)  e  \textbf{Rússia} (LOC),  já  haviam  escolhido  os  seus representantes para esse mesmo ano. Conforme o desenrolar da situação, e pela tomada de  posição  dos  diferentes  órgãos  e  países,  também  o  \textbf{Comité  Executivo} (ORG)  decidiu  tomar posição  contrariando  a  sua  neutralidade,  e  excluir  a  participação  russa  da  competição.

Para  um  festival  que  se  pretende  mostrar  imparcial  e  apolítico,  esta  foi  uma  decisão parcial,  ao  tomar  partido  por  um  dos  lados,  seguindo  o  modelo  de  outras  instituições europeias, no sentido de sancionar o país russo.

Também na edição de 2023, conseguimos encontrar algumas referências, seja na letra  ou  na performance de  um  determinado  país,  sobre  este  acontecimento.  \textbf{Vesna} (MISC), representantes da \textbf{República Checa} (LOC), apresentaram-se em palco com a canção \textbf{My Sister's Crown} (MISC) ,  cuja  letra  cantada  nas  línguas  inglesa,  ucraniana,  checa  e  búlgara  nos  fala  da resistência de uma "irmã", que se mantém forte e verdadeira a si mesma, promovendo a paz e o amor, em vez da guerra e de perdas humanas: \textbf{My sister won't stand in the corner/} (MISC) Nor will she listen to you/ \textbf{My sister} (MISC), \textbf{wild at heart/ Will} (MISC) never let you tie her down (…) \textbf{My beautiful sister} (MISC), you are so \textbf{strong/ Brave and the only one} (MISC), \textbf{the crown is yours} (MISC) (…) \textbf{Give me your hand don't be afraid/} (MISC) \textbf{Come swim} (MISC) with your sisters/ There's no place for hate in \textbf{our sea} (ORG) (…) You can take your hands back/ No \textbf{one wants more boys dead} (MISC).

Ao  lado  desta  composição,  encontramos  também  os  representantes  da  \textbf{Croácia} (LOC), \textbf{Let 3} (MISC), com a música \textbf{Mama} (MISC) ŠČ!

, cantada na língua croata, na qual é possível encontrar, na  letra,  semelhanças  com  a  canção  anterior.  Na  atuação  da  \textbf{Croácia} (LOC),  é  possível depreender-se uma posição negativa perante o referido confronto bélico, sobretudo pela forma como se apresentaram em palco. "\textbf{Mama} (MISC), \textbf{I'm} (ORG) going to war".

Em  anos  anteriores,  também  se  encontram  semelhanças  no  que  diz  respeito  a situações  e,  por  conseguinte,  posições  políticas. \textbf{Relembremos} (ORG)  o  conflito entre  \textbf{Israel} (LOC)  e \textbf{Palestina} (LOC), considerado um dos mais longos da \textbf{História} (MISC). Numa situação em que apenas um dos integrantes faz parte do \textbf{FEC} (ORG), as tomadas de posição poderiam  não ser tão visíveis como no caso da \textbf{Ucrânia} (LOC). Porém, em plena \textbf{Green Room} (ORG) , os representantes da \textbf{Islândia} (LOC), em  2019,  \textbf{Hataria} (LOC),  tomaram  a  posição  do  lado  palestiniano  quando,  no  momento  de distribuição dos votos, ergueram uma bandeira da \textbf{Palestina} (LOC).

https://www.youtube.com/watch?v=ag8qxpvTTy0 - \textbf{Vesna} (MISC): \textbf{My Sister's Crown} (MISC). (A minha irmã não vai ficar no canto/ \textbf{Nem te} (MISC) vai ouvir/ A minha irmã, selvagem de coração/ Nunca vai deixar que a amarres (...) Minha linda irmã, és tão forte/ \textbf{Corajosa} (LOC) e a única, a coroa é tua (...) Dá-me a tua mão, não tenhas medo/ \textbf{Vem} (PER) nadar com as tuas irmãs/ Não há lugar para o ódio no nosso mar (...) \textbf{Podes} (LOC) recuperar as tuas mãos/ Ninguém quer mais rapazes mortos). (\textbf{Tradução nossa} (MISC)) - No anexo 3 encontramos a letra completa.

https://www.youtube.com/watch?v=hGuGfdEJ5Pw - \textbf{Let 3} (MISC): \textbf{Mama} (MISC) ŠČ!

. (\textbf{Mamã} (LOC), eu vou para a guerra).

(\textbf{Tradução nossa} (MISC)) - No anexo 4 encontramos a letra completa.

\textbf{Aviv} (PER).

Esse  mesmo  ato  valeu-lhes  uma  multa,  uma  vez  que  "\textbf{Contest  organizers  the European Broadcasting Union} (MISC) said \textbf{Saturday in} (MISC) a statement sent to \textbf{AFP that the} (MISC) gesture \textbf{infringed} (PER) their rules banning political gestores\textbf{"} (MISC) (\textbf{AFP} (MISC), 2019: s/p).

Para além dos já apontados, muitos são os momentos  e manifestações políticas que poderíamos elencar como boicotes: due to voting or political \textbf{situation in} (MISC) a host state (\textbf{Spain} (PER)), false broadcasting in \textbf{Jordan in order to hide the victory of Israel} (ORG) is a reflection of tangled bilateral relations, as well as \textbf{expulsion from} (MISC) the ESC due to the unpaid debt to \textbf{EBU} (MISC) (\textbf{Romania} (LOC)) (\textbf{Borić} (LOC), 2017: 228).

Em anos anteriores, outros países que antes pertenciam à denominada \textbf{Jugoslávia} (ORG), tais como \textbf{Bósnia e Herzegovina} (LOC), \textbf{Croácia} (LOC) e \textbf{Eslovénia} (LOC), assumiram posições referentes a conflitos políticos: \textbf{One of the most striking examples of using the ESC} (MISC) for sending a political \textbf{message by using the soft} (MISC), nonpolitical, medium in 1993 was \textbf{participation of Bosnia and Herzegovina} (MISC) (\textbf{Jordan} (ORG), 2014). \textbf{Since at that time the war was raging in  former} (MISC) \textbf{Yugoslavia} (LOC), particular \textbf{attention was given to performers from} (MISC) \textbf{Bosnia and Herzegovina} (LOC), \textbf{Croatia} (LOC), and \textbf{Slovenia} (LOC), the  newly  independent  countries  which  made  their  debut  \textbf{in  the  competition  after  the dissolution of Yugoslavia} (MISC). \textbf{Since} (PER) the global audience at that time was not getting enough \textbf{information about the war on the European} (MISC) soil (…) \textbf{Their} (LOC) entry songs strongly referred to the turmoil in the region which, in 1993, showed no \textbf{signs of slowing down} (MISC). Croatia participated with a song titled "\textbf{Don't ever cry} (MISC)" \textbf{which talked} (PER) about a young man Ivan who died in the war, and \textbf{Bosnian} (MISC) entry song was \textbf{"Sva} (LOC) bol svijeta" (\textbf{All the Pain in the World} (MISC)) \textbf{which had} (PER) a similar undertone (\textbf{Borić} (LOC), 2017: 229/230).

A tomada de posição a nível político no cerne do \textbf{FEC} (ORG) irá estar sempre presente, seja de forma percetível, seja de forma mais dissimulada. Tal poderá ser explicado devido à convivência multicultural do Velho Continente assim como à resistência e à natureza reformadora de todos os países integrantes.

Uma  tendência  que  já  se  começou  a  impor  no  \textbf{FEC} (ORG)  é,  como  já  abordámos anteriormente,  as  redes  sociais.  Com  a  criação  deste  novo  modo  de  comunicar  e  o posicionamento do concurso e dos países cada mais se vive um caminho pré-festival que liga o espetador ao artista. Mais recentemente, é possível perceber que nas plataformas digitais  a  opinião  pública  pode  contribuir  para  ditar  o  resultado  final  da  competição.

Contudo,  o  fenómeno  reverso  pode  também  acontecer.  Na  edição  de  2023  do  \textbf{FEC} (MISC),  o grande  vencedor  do  público  era  \textbf{Käärijä} (PER),  com \textbf{Cha  Cha  Cha} (MISC) ,  assim  como  também  a representante norueguesa, \textbf{Alessandra} (PER). Se recordarmos a edição 2019, o mesmo  se terá passado com \textbf{Mahmood} (LOC), representante da \textbf{Itália} (LOC). As redes sociais, com maior incidência no \textbf{Twitter} (MISC) e no \textbf{Tiktok} (PER) , permitem uma prospeção e a colocação do artista no centro das conversas sobre o \textbf{FEC} (ORG). Inclusive, podemos encontrar composições que normalmente não conseguiriam lugar numa final, mas que se conseguem apurar por conter uma determinada parte  que  cativou  o  público,  como  é  o  caso  da  música \textbf{Solo} (MISC) ,  da  \textbf{Polónia} (LOC),  em  2023.  A primeira palavra, respetivamente baby , com o ritmo da música e da entoação da artista tornaram-se virais.

O pós-final eurovisivo conta também com a exposição e visibilidade dos artistas, não só pela sua vitória, mas também pelos seus looks, como foi o caso dos vencedores de 2021, os \textbf{Måneskin} (LOC). Este grupo conquistou a \textbf{Europa} (LOC) e o mundo de rompante, levando a que o pop rock voltasse a estar no centro da ação, como no século passado: \textbf{But  Måneskin} (LOC),  whose  name  \textbf{is  the  Danish  word  for  moonlight} (MISC),  has  instead  achieved  a level of global popularity \textbf{unprecedented in the history of Italian} (MISC) rock. And \textbf{they've} (ORG) done it faster than any \textbf{Eurovision winner ever} (ORG), another indication of how quickly stardom can happen in the digital age (\textbf{Shaw} (PER), 2021: s/p).

Este  crescimento  foi  largamente  estimulado  pelo \textbf{Tiktok} (PER) ,  cujo  algoritmo  pode favorecer um determinado vídeo e levá-lo a milhões de visualizações. Na edição de 2022, embora não tenha conseguido um lugar no top 10 , \textbf{Rosa Linn} (PER), representante da \textbf{Arménia} (LOC), viu  a  sua  canção  eurovisiva  tornar-se  num hit mundial  nesta  mesma  plataforma.  Este algoritmo permitiu que a música Snap conquistasse de rompante o público mais afastado da \textbf{Eurovisão} (LOC). "\textbf{Rosa Linn's Snap} (PER) has been used on more than 360,000 \textbf{TikTok} (MISC) clips, \textbf{with some having} (PER) millions of likes (…) \&\#34;You never know what will go viral on \textbf{TikTok,\&\#34} (LOC); \textbf{Rosa Linn} (PER) said" (\textbf{BBC} (ORG), 2022: s/p).

O  \textbf{FEC} (MISC)  está  sempre  em  constante  mudança  e  adaptação  aos  novos  meios  e  à sociedade  em  que  se  insere.  Todos  os  desafios,  oposições  e  tomadas  de  posição demonstram o poder de união através de um movimento cultural e como este pode unir nas diferenças. É sabido que, contendo as suas regras, os artistas, caso queiram apresentar e motivar algum assunto ou finalidade, devem de o fazer de forma a não prejudicar ou injuriar outro concorrente.

(…) I intentionally draw attention to the cultural and performance sites \textbf{at which the Eurovision Song Contest is open for the emergence of many different voices of European} (MISC) difference.

\textbf{Philip Vilas Bohlman} (PER) 3.

O \textbf{Festival Eurovisão da Canção} (MISC) para a comunidade europeia Como  tentámos  demonstrar  ao  longo  dos  dois  pontos  anteriores,  existe  uma ligação  intrínseca  entre  o  \textbf{FEC} (MISC)  e  a  \textbf{comunidade  europeia} (ORG).  Não  sendo  exclusiva,  essa conexão pode facilitar um caminho comunicacional entre ambas as partes. O facto de ser produzido pela \textbf{EBU} (MISC) fundamenta essa relação e torna a competição mais específica para o seu raio de influência. "\textbf{The possibility that transnational TV} (MISC) programes such as \textbf{ESC} (MISC) might  contribute  to  a  sense  \textbf{of} (MISC)\textbf{European  belonging  has  been  a} (MISC)\textbf{point  of  some  debate} (LOC)" (\textbf{Pajala} (LOC), 2013: 79). Atualmente, o seu raio de influência ultrapassa as fronteiras europeias, o que permite demonstrar a sua expansão no seio da sociedade e da comunidade em que o mesmo se desenrola: A  closer  look  \textbf{at  the} (MISC)  data  shows  that  cases  in  which  the  audiovisual  construction  adds \textbf{Europeanness} (ORG)  to  the  lyrics  hardly  ever  occur,  that  is,  when  transnational  \textbf{European construction} (MISC) takes place, \textbf{the linguistic level is generally involved} (ORG) (\textbf{mostly by choice  of non-national European languages} (MISC)) (\textbf{Motschenbacher} (MISC), 2016: 251).

Trata-se de um festival que surge como um elemento unificador de países e, assim também de diferentes culturas, que tem vindo a aglutinar cada vez mais países ao longo das suas já muitas edições. Até à data, integraram pela primeira vez o \textbf{FEC} (ORG) as seguintes nações: - 1956: \textbf{Bélgica} (LOC), \textbf{França} (LOC), \textbf{Alemanha} (LOC), \textbf{Itália} (LOC), \textbf{Luxemburgo} (LOC), \textbf{Países Baixos} (LOC) e \textbf{Suíça} (LOC); - 1957: \textbf{Áustria} (LOC), \textbf{Reino Unido} (LOC) e \textbf{Dinamarca} (LOC); - 1959: \textbf{Mónaco} (LOC); - 1960: \textbf{Noruega} (LOC); - 1961: \textbf{Espanha} (LOC), \textbf{Finlândia} (LOC), \textbf{Jugoslávia} (ORG) e \textbf{Suécia} (LOC); - 1964: \textbf{Portugal} (LOC); - 1965: \textbf{Irlanda} (LOC); - 1971: \textbf{Malta} (LOC); - 1973: \textbf{Israel} (LOC); - 1974: \textbf{Grécia} (LOC); - 1975: \textbf{Turquia} (LOC); - 1980: \textbf{Marrocos} (LOC); - 1981: \textbf{Chipre} (LOC); - 1986: \textbf{Islândia} (LOC); - 1993: \textbf{Eslovénia} (LOC), \textbf{Bósnia e Herzegovina} (LOC) e \textbf{Croácia} (LOC); - 1994: \textbf{Estónia} (LOC), \textbf{Lituânia} (LOC), \textbf{Polónia} (LOC), \textbf{Hungria} (LOC), \textbf{Rússia} (LOC), \textbf{Eslováquia} (LOC) e \textbf{Roménia} (LOC); - 1996: \textbf{Macedónia do Norte} (LOC); - 2001: \textbf{Letónia} (LOC); - 2003: \textbf{Ucrânia} (LOC); - 2004: \textbf{Sérvia e Montenegro} (LOC), \textbf{Albânia} (LOC), \textbf{Andorra} (LOC) e \textbf{Bielorrússia} (LOC); - 2005: \textbf{Moldávia} (LOC) e \textbf{Bulgária} (LOC); - 2006: \textbf{Arménia} (LOC); - 2007: \textbf{Montenegro} (LOC) e \textbf{República Checa} (LOC); - 2008: \textbf{São Marino} (LOC), \textbf{Geórgia} (LOC) e \textbf{Azerbaijão} (LOC); - 2010: \textbf{Sérvia} (LOC).

Ao lado da comunidade europeia, enquanto participante, estiveram, estão e podem vir a estar vários países periféricos do continente europeu. Exemplo disso poderá ser a participação da \textbf{Austrália} (LOC) que surgiu, num primeiro momento, como complementaridade.

Esta competição funciona como um caminho para a compreensão e para o diálogo, uma união para observar uma manifestação cultural que pode conectar pessoas de diferentes culturas, idades e passados. Desde o início do \textbf{FEC} (ORG), existe em determinados países um convívio quando este acontecimento ocorre. As famílias e os grupos de amigos juntam- se para celebrar esta competição proporcionada pela \textbf{EBU} (MISC). Com o desenvolvimento das redes  sociais,  o fandom eurovisivo  tem  crescido  exponencialmente  e  tem  estado  mais próximo daqueles que, em períodos anteriores, estavam mais afastados.

Um pouco por toda a parte, existem, para além das transmissões no canal oficial da  \textbf{Eurovisão} (LOC)  e  no \textbf{Youtube} (MISC) , viewing  parties que  permitem  transmitir  os  seus eventos.  Estas têm o poder de unir os fãs eurovisivos e de, através da reunião à volta da televisão ou de uma plataforma digital, demonstrar o seu apoio, num raio mais pequeno de  influência.  "\textbf{Gatherings} (MISC)  or  parties  are  importante  as  they  secure  \textbf{Eurovision's} (MISC)  place within familial or social rituals. What occurs at these parties helps to create new memories for that group\textbf{"} (MISC) (\textbf{Carniel} (MISC), 2018: 53). O que significa que o \textbf{FEC} (ORG) sempre foi visto como um elemento de união, seja familiar, político, ou de amizade. Essa sua natureza permite-lhe cativar e arrecadar cada vez mais reconhecimento intra e extraeuropeu porque, apesar de ocorrer num contexto mais pequeno, promove o debate sobre a diversidade, a diferença e a aceitação do outro.

3.1.

\textbf{Unidade} (LOC) e \textbf{Diversidade} (LOC) cultural (vista a partir de um inquérito) Para  a  comunidade  europeia,  o  \textbf{FEC} (MISC)  detém  um  papel  relevante.  Mas,  para  que fosse  possível  uma  tomada  de  conhecimento  fidedigna,  sobre  uma  amostra  razoável, colocámos em prática um inquérito com o objetivo de saber se existe um sentimento de unidade  ou  diversidade  no  seio  da  \textbf{Eurovisão} (LOC).  O  inquérito  esteve  disponível  para preenchimento entre os dias 3 de fevereiro a 3 de abril de 2023, através das redes sociais, \textbf{WhatsApp} (MISC) , \textbf{Facebook} (MISC) e \textbf{Instagram} (ORG) .

A  amostra  contou  com  um  total  de  95  respostas.  Quanto  a  variáveis sociodemográficas, obtivemos os seguintes resultados: - 80\% do sexo feminino; - 15,8\% do sexo masculino; - 4,2\% preferiram não responder.

No  que  diz  respeito  à  faixa  etária,  a  que  se  destaca  com  um  maior  número  de respostas é a dos maiores de 50 anos, com 42,1\%, seguida da faixa etária entre 40 e 49 anos, com 25,3\%. As idades referentes a menos de 20 anos e entre 30 e 39 anos não foram além dos 3,2\% e dos 6,3\%, respetivamente. Quanto à área de residência dos inquiridos, os  valores  mais  elevados  registados  foram  \textbf{Funchal} (LOC)  com,  52,6\%,  \textbf{Caniço} (MISC),  11,6\%, \textbf{Camacha} (LOC),  8,4\%  e  \textbf{Santa  Cruz} (LOC),  6,3\%,  todos  da  \textbf{Região  Autónoma  da  Madeira} (LOC)  (\textbf{RAM} (MISC)).

Também obtivemos respostas de locais fora da \textbf{RAM} (MISC), tais como \textbf{Jersey} (LOC), \textbf{Lisboa} (LOC), \textbf{Amadora} (LOC), \textbf{Austrália} (LOC),  \textbf{Mafra} (LOC)  e  \textbf{Porto} (LOC),  num  total  de  9,7\%.  Em  termos  de  habilitações  académicas, destacam-se  43,2\%  dos  inquiridos  com  licenciatura,  23,2\%  com  ensino  secundário  e 11,6\%  com  3  ciclo/  9  ano.  As  percentagens  mais  baixas  registadas  referem-se  ao  1 ciclo/  ensino  primário,  bacharelato  e  doutoramento,  com  um  total  de  2,1\%.  No  que concerne a ocupação dos mesmos, 77,9\% encontra-se empregado, 12,6\% são estudantes e 9,5\% estão desempregados.

Todos os inquiridos afirmaram conhecer o \textbf{FEC} (ORG). Todavia, 52,6\% acompanham às vezes,  25,3\%  acompanham  sempre,  17,9\%  acompanham  raramente  e  4,2\%  nunca acompanham.  No  que  diz  respeito  à  pergunta  sobre  as  razões  que  levam  a  assistir  o festival, as respostas indicaram os seguintes resultados: - 45,1\% para apoiar o seu país numa competição europeia; - 31,9\% gostam de observar diferentes aspetos culturais; - 31,9\% acham interessante; - 27,5\% por curiosidade; - 1,1\% está ligado ao mundo da música; - 1,1\% porque há décadas que não há um festival que interesse; - 1,1\% para avaliar a qualidade musical; - 1,1\% acompanhar artistas de que gosta; - 1,1\% algumas tendências ou os \textbf{Media} (LOC) promovem a \textbf{Eurovisão} (LOC) e devido a participações específicas que tiveram boas repercussões no público.

Sobre  o  impacto  que  o  \textbf{FEC} (MISC)  pode  ter  na  população  e  no  contexto  intra  e extraeuropeu, obtivemos as seguintes respostas: - 64,2\% acreditam que promove a representatividade de diferentes gentes e culturas; - 27,4\% alegam que é um movimento cultural; - 26,3\% afirmam que pode alterar a visibilidade de um determinado país; - 1,1\% encontra corrupção nos direitos televisivos e influencia nos resultados; -  1,1\%  acha  que  promove  a  dedicação  e  valorização  do  património  local  com  as  suas adaptações.

Quanto à pergunta se, desde a fundação, o seu país já tinha vencido: - 86,3\% atestam que sim; - 9,5\% referem não ter esse conhecimento.

É certo que, para que possa alterar a visibilidade do país, uma vitória pode facilitar esse mesmo feito. Dos inquiridos, - 87,1\% atestam que este facto pode ocorrer; - 10,6\% não o sabem; - 2,4\% acreditam que a vitória não altera a sua visibilidade.

Sendo o nosso objeto de estudo o \textbf{FEC} (ORG) e as manifestações culturais de cada país dentro  do  mesmo,  os  elementos  culturais  que  já  tivemos  oportunidade  de  abordar  em capítulos  anteriores,  como  a  língua  e  o  vestuário,  podem  destacar-se  como  motor diferencial entre os mesmos. De acordo com 86,3\% dos respondentes, aspetos como estes podem conduzir a essa diferenciação, 12,6\% afirmam que talvez o facto aconteça. Estes aspetos distintos, consoante os inquiridos, podem ser representativos de uma cultura e de um povo, tendo sido declarado por 83,2\% dos participantes. Por sua vez, 12,6\% afirmam que talvez a língua, as vestes e as melodias sejam importantes para a representatividade de uma cultura e 4,2\% não o consideram.

Devido  à  sua  natureza  e  às  proporções  que  o  \textbf{FEC} (MISC)  detém,  também  podemos encontrar  diversas  vantagens  no  mesmo.  Algumas  das  respostas  mais  comuns  dos inquiridos ressalvam: - A visibilidade e promoção que poderia vir a ter um país; - Dar a conhecer diferentes traços culturais de diferentes povos europeus; - Aumento do turismo devido à divulgação da nação; - Promoção de novos talentos e da diversidade cultural.

Porém, foram apontadas as seguintes desvantagens: - A utilização e uniformização de uma determinada língua (o inglês); - \textbf{Interesses} (LOC) económicos; - Produção de estereótipos; - O sistema de voto e imparcialidade.

Estas questões foram apresentadas aos inquiridos como "resposta aberta" para que pudéssemos obter pontos de vista distintos, que pudessem contribuir e valorizar a nossa investigação,  uma  vez  que  procuramos  conhecer  o  sentimento  das  pessoas  de  unidade e/ou diversidade dentro do contexto europeu, através do \textbf{FEC} (ORG).

Com os resultados advindos do inquérito que utilizámos é possível tirar algumas conclusões  que  nos  parecem  pertinentes  para  a  nossa  investigação.  Assim,  podemos concluir que o \textbf{FEC} (ORG) está muito presente na vida dos europeus, mesmos que estes não o acompanhem  regularmente.  Esta  competição  move  as  pessoas  de  uma  maneira semelhante às demais competições, podendo estas estar ligadas ao mundo do desporto, por  exemplo,  no  sentido  de  querer  apoiar  a  sua  nação  e  ver-se  representado  além- fronteiras,  ao  mais  alto  nível  possível.  A  convergência  e  convivência  dos  diferentes países, gentes e  culturas  favorece um melting-pot dentro do evento,  criando condições para um prolongamento da ligação mesmo após terminar. O evento promove junto dos povos, assim, a representatividade cultural e a multiculturalidade de uma forma simples e  abrangente.  A  língua,  o  vestuário,  as  melodias,  entre  outros  elementos  que  podem culturalmente  distinguir  os  indivíduos  entre  si,  são  valorizados  pelo  público  porque permitem  exaltar  uma  determinada  cultura.  Essa  particularidade  confere  ao  \textbf{FEC} (MISC)  a unidade através da diversidade e não o sentimento de um ou de outro. É claro que, nem sempre conseguimos encontrar este sentimento de distinção e podemos vir a encontrar outras situações. A homogeneização de uma língua, especialmente nos últimos anos, tem vindo  a  tornar-se  algo  comum  o  que  pode  retirar,  de  certa  forma,  a  individualidade  e identificação de uma nação  consoante  a língua em  que se manifesta. Pode existir uma descaracterização cultural por ser adotada uma língua e estilos musicais distintos daqueles que são inerentes à sua cultura, encontrando uma uniformização linguística e musical.

Também, direta ou indiretamente, apesar de o \textbf{FEC} (ORG)  ser considerado um festival apolítico,  estas  questões  podem  vir  à  tona  e  criar  situações  de  instabilidade  que  nada teriam a ver com a competição. Normalmente, tratando-se de um concurso em que reina a meritocracia, situações como a da \textbf{Ucrânia} (LOC) em anos mais recentes podem fomentar uma vitória politicamente correta, retirando o mérito aos seus adversários. Inclusive, o sistema de  voto  pode  estar  a  desviar-se  das  músicas  favoritas  e  da  sua  pontuação,  passando  a deter-se sobre outras questões, tais como a situação do país, a intervenção dos artistas, entre  outros.  Em  vez  de  ser  representado  pelos  seus  hábitos,  raízes  e  costumes,  o  país pode  ser  associado  a  interesses  e  opiniões  políticas,  apesar  de  o  mesmo  ser  feito  com referências subtis. Com as crescentes minorias no território europeu, estejam estas ligadas à etnia, sexualidade, religião, entre outras, começam a tornar-se cada vez mais visíveis em palco, apesar de \textbf{"The Eurovision Song Contest} (MISC) has never been about minorities and minority identities, in general or specific" (\textbf{Bohlman} (PER), 2007: 47). Essa chamada permite que  exista  uma  discussão  sobre  a  aceitação  das  diferenças  e  promove  a  \textbf{união} (LOC)  pela diversidade, podendo vir a criar uma maior tolerância ao diferente.

Mas  não  só  de  pontos  contrários  se  sustenta  o  \textbf{FEC} (MISC).  A  opinião  da  maioria  dos inquiridos parece coincidir em algumas das ideias que resultam da nossa investigação. O \textbf{FEC} (ORG) funciona como motor para a diversidade cultural dentro do \textbf{Velho Continente} (LOC), mas é nessa diversidade que reside a unidade que os indivíduos podem vir a sentir. Os povos podem  encontrar-se  refletidos,  através  dos  seus  semelhantes,  criando  um  ambiente  de cooperação e proximidade. Não só para os fatores individualistas, mas quão mais refletido esteja  uma  nação,  mais  repercussões  positivas  podem  daí  advir.  O  turismo  é  um  dos movimentos  que  mais  se  expande  e  é  cada  vez  mais  favorecido  pela  competição.  Ao colocarmos  uma  cultura  no  centro  das  atenções,  conseguimos  retirar  conclusões  sobre diversos pontos que podem vir a propiciar aos participantes um conhecimento mais geral sobre o que é aquela nação: \textbf{Fan  tourism  traditionally  often} (LOC)  takes  the  form  of  a  secular  pilgrimage  (\textbf{Hall} (LOC),  2002; \textbf{Digance} (MISC), 2006), as a location or site that has a meaning in the "text" surrounding a popular culture figure may be regarded as \textbf{"} (ORG)sacred" \textbf{within the fandom} (MISC) (\textbf{Linden} (LOC) \&amp; \textbf{Linden} (LOC), 2017).

(…) \textbf{It is clear that fans and fandoms} (MISC) play a key role in popular culture travel - and in particular in event-led travel (\textbf{Linden} (LOC), 2018: 1/11).

Este inquérito permitiu-nos perceber que, apesar de não ser um tema recorrente ao longo do ano, o \textbf{FEC} (ORG) consegue despertar a atenção dos telespectadores na altura em que este ocorre, falando do público em geral. O seu fandom acompanha com  mais atenção todos os passos até à celebração final do evento. Para todos os inquiridos, o \textbf{FEC} (MISC) significa algo,  mesmo  que  o  seu  conhecimento  sobre  a  competição  não  seja  exaustivo.  Neste sentido, é possível afirmar que o festival faz parte de grande parte da comunidade em que nos encontramos inseridos.

\textbf{Conclusão Como} (MISC)  tentámos  demonstrar  ao  longo  da  nossa  investigação,  o  sentimento  de unidade  através  da  diferença  complementam-nos  e  tendem  a  estar  na  essência  do continente  europeu  e,  em  especial,  no  \textbf{Festival} (MISC)Eurovisão  da  Canção.  Os  elementos históricos que marcaram o Velho Continente fundamentam as relações intracontinentais e extracontinentais, que se encontram na base de qualquer forma de comunicação entre os  povos.  A  língua,  através  das  suas  raízes  etimológicas,  permite  um  senso  de compreensão e abertura ao semelhante. Também as tradições que se assemelham entre si no espaço de influência permitem uma adaptação mais simplificada do indivíduo quando este passa a integrar uma cultura distinta daquela que o viu nascer. Da nossa investigação cremos poder afirmar que o \textbf{Festival Eurovisão da Canção} (MISC) coloca em palco a cultura, os elementos culturais e o bom nome de cada nação que integra esse encontro musical e de diferentes povos, concedendo a estes um espaço neutro em que a sua substância e essência pode ser observada pelos demais.

É através da conclusão deste percurso histórico e cultural europeu que julgamos poder  afirmar  que  existe  um  sentimento  de  unidade  e  diversidade  cultural  no  espaço europeu. As influências de várias partes do continente e do globo, através das migrações, favorecem também a multiculturalidade e, apesar de em algumas situações encontrarmos alguns entraves à convivência das gentes, a sociedade do século \textbf{XXI} (MISC) parece encontrar-se mais apta e mais aberta para a comunicação e compreensão  da diversidade. O \textbf{Festival Eurovisão da Canção} (MISC) propicia este acontecimento desde a sua fundação porque sempre procurou promover os diferentes passados e os aspetos culturais distintos de todos aqueles que o integram, demonstrando a premissa de aceitação do diferente, do estrangeiro, do diverso.

Este estudo permite-nos ajudar a perceber como uma competição musical europeia consegue criar no espectador diferentes repercussões, seja a nível do turismo, do gosto pela  cultura,  entre  outros.  \textbf{Procurámos} (MISC),  assim,  demonstrar  como  esta  manifestação cultural, reconhecível além-fronteiras, se torna numa fonte de inspiração para modelos semelhantes, tais como o \textbf{American Song Contest} (ORG) ou o \textbf{Eurovision Young Musicians} (MISC). A opinião  pública  manifestada  através  dos  inquéritos,  realizados  via  online,  veio  trazer respostas a várias das interrogações que colocámos à partida para a nossa investigação.

Através dos inquiridos julgamos poder comprovar que o \textbf{FEC} (ORG) continua a ser um ponto de encontro e de partilha de diversos povos e culturas, promovendo a sua difusão ao mundo.

Não  só,  mas  também  as  línguas  nacionais  utilizadas,  apesar  de  terem  perdido  espaço, quando apresentadas em palco e em competição arrecadam bons lugares, favorecendo a singularidade que cada uma destas promove face o \textbf{Festival Eurovisão da Canção} (MISC).



\vspace{1cm}
\begin{thebibliography}{9}
\bibitem{content}
    Fonte analisada: \texttt{content}. Texto processado automaticamente.
\end{thebibliography}

\end{document}
